
\Data\ID{1003019901}
\Updated{2022-08-25 10:08:02}
\Level{A}
\SubArea{Radical Equations and Inequalities}
\LANGen{
\Question{
Find the solution of the equation.
\[
\sqrt{-3x^2+2x+5}=3
\]}
{\(x\in\emptyset\)}{\(x\in[-2;2]\)}{\(x\in\{-2;2\}\)}{\(x\in\mathbb{R}\)}{}{}}
\LANGpl{
\Question{
Rozwiąż równanie.
\[
\sqrt{-3x^2+2x+5}=3
\]}
{\(x\in\emptyset\)}{\(x\in\langle-2;2\rangle\)}{\(x\in\{-2;2\}\)}{\(x\in\mathbb{R}\)}{}{}}
\LANGcs{
\Question{
Určete řešení rovnice.
\[
\sqrt{-3x^2+2x+5}=3
\]}
{\(x\in\emptyset\)}{\(x\in\langle-2;2\rangle\)}{\(x\in\{-2;2\}\)}{\(x\in\mathbb{R}\)}{}{}}
\LANGsk{
\Question{
Vyriešte rovnicu.
\[
\sqrt{-3x^2+2x+5}=3
\]}
{\(x\in\emptyset\)}{\(x\in\langle-2;2\rangle\)}{\(x\in\{-2;2\}\)}{\(x\in\mathbb{R}\)}{}{}}
\LANGes{
\Question{
Halla la solución de la siguiente ecuación.
\[
\sqrt{-3x^2+2x+5}=3
\]}
{\(x\in\emptyset\)}{\(x\in[-2;2]\)}{\(x\in\{-2;2\}\)}{\(x\in\mathbb{R}\)}{}{}}
\End

\Data\ID{1003019902}
\Updated{2022-08-25 10:08:02}
\Level{A}
\SubArea{Radical Equations and Inequalities}
\LANGen{
\Question{
Find the solution of the equation.
\[
\sqrt{x^2-4x+3}=\sqrt{x^2-6x+3}
\]}
{\(x\in\{0\}\)}{\(x\in\mathbb{R}\)}{\(x\in\emptyset\)}{\(x\in\mathbb{R}\setminus\{0\}\)}{}{}}
\LANGpl{
\Question{
Rozwiąż równanie.
\[
\sqrt{x^2-4x+3}=\sqrt{x^2-6x+3}
\]}
{\(x\in\{0\}\)}{\(x\in\mathbb{R}\)}{\(x\in\emptyset\)}{\(x\in\mathbb{R}\setminus\{0\}\)}{}{}}
\LANGcs{
\Question{
Určete řešení rovnice.
\[
\sqrt{x^2-4x+3}=\sqrt{x^2-6x+3}
\]}
{\(x\in\{0\}\)}{\(x\in\mathbb{R}\)}{\(x\in\emptyset\)}{\(x\in\mathbb{R}\setminus\{0\}\)}{}{}}
\LANGsk{
\Question{
Vyriešte rovnicu.
\[
\sqrt{x^2-4x+3}=\sqrt{x^2-6x+3}
\]}
{\(x\in\{0\}\)}{\(x\in\mathbb{R}\)}{\(x\in\emptyset\)}{\(x\in\mathbb{R}\setminus\{0\}\)}{}{}}
\LANGes{
\Question{
Halla la solución de la siguiente ecuación.
\[
\sqrt{x^2-4x+3}=\sqrt{x^2-6x+3}
\]}
{\(x\in\{0\}\)}{\(x\in\mathbb{R}\)}{\(x\in\emptyset\)}{\(x\in\mathbb{R}\setminus\{0\}\)}{}{}}
\End

\Data\ID{1003020601}
\Updated{2022-08-25 10:08:13}
\Level{A}
\SubArea{Radical Equations and Inequalities}
\LANGen{
\Question{
Find the domain of the expression.
\[
\frac1{\sqrt{5x^2+7x-6}}
\]}
{\(\left(-\infty;-2\right)\cup\left(\frac35;\infty\right)\)}{\(\left(-2;\frac35\right)\)}{\(\left(-\infty;-2\right]\cup\left[\frac35;\infty\right)\)}{\(\left[-2;\frac35\right]\)}{}{}}
\LANGpl{
\Question{
Znajdź dziedzinę podanego wyrażenia.
\[
\frac1{\sqrt{5x^2+7x-6}}
\]}
{\(x\in\left(-\infty;-2\right)\cup\left(\frac35;\infty\right)\)}{\(x\in\left(-2;\frac35\right)\)}{\(x\in\left(-\infty;-2\right\rangle\cup\left\langle\frac35;\infty\right)\)}{\(x\in\left\langle-2;\frac35\right\rangle\)}{}{}}
\LANGcs{
\Question{
Určete definiční obor výrazu.
\[
\frac1{\sqrt{5x^2+7x-6}}
\]}
{\(\left(-\infty;-2\right)\cup\left(\frac35;\infty\right)\)}{\(\left(-2;\frac35\right)\)}{\(\left(-\infty;-2\right\rangle\cup\left\langle\frac35;\infty\right)\)}{\(\left\langle-2;\frac35\right\rangle\)}{}{}}
\LANGsk{
\Question{
Určte definičný obor výrazu.
\[
\frac1{\sqrt{5x^2+7x-6}}
\]}
{\(\left(-\infty;-2\right)\cup\left(\frac35;\infty\right)\)}{\(\left(-2;\frac35\right)\)}{\(\left(-\infty;-2\right\rangle\cup\left\langle\frac35;\infty\right)\)}{\(\left\langle-2;\frac35\right\rangle\)}{}{}}
\LANGes{
\Question{
Halla el dominio de la siguiente expresión.
\[
\frac1{\sqrt{5x^2+7x-6}}
\]}
{\(\left(-\infty;-2\right)\cup\left(\frac35;\infty\right)\)}{\(\left(-2;\frac35\right)\)}{\(\left(-\infty;-2\right]\cup\left[\frac35;\infty\right)\)}{\(\left[-2;\frac35\right]\)}{}{}}
\End

\Data\ID{1003020602}
\Updated{2022-08-25 10:08:13}
\Level{A}
\SubArea{Radical Equations and Inequalities}
\LANGen{
\Question{
Find the domain of the expression.
\[
\frac{\sqrt{x^2-x-6}}{x^2-7x+10}
\]}
{\(\left(-\infty;-2\right]\cup\left[3;5\right)\cup\left(5;\infty\right)\)}{\(\left(-\infty;-2\right)\cup\left(3;\infty\right)\)}{\(\left(-\infty;-2\right)\cup\left(3;5\right)\cup\left(5;\infty\right)\)}{\(\left(-\infty;2\right)\cup\left(3;5\right)\cup\left(5;\infty\right)\)}{}{}}
\LANGpl{
\Question{
Znajdź dziedzinę podanego wyrażenia.
\[
\frac{\sqrt{x^2-x-6}}{x^2-7x+10}
\]}
{\(x\in\left(-\infty;-2\right\rangle\cup\left\langle3;5\right)\cup\left(5;\infty\right)\)}{\(x\in\left(-\infty;-2\right)\cup\left(3;\infty\right)\)}{\(x\in\left(-\infty;-2\right)\cup\left(3;5\right)\cup\left(5;\infty\right)\)}{\(x\in\left(-\infty;2\right)\cup\left(3;5\right)\cup\left(5;\infty\right)\)}{}{}}
\LANGcs{
\Question{
Určete definiční obor výrazu.
\[
\frac{\sqrt{x^2-x-6}}{x^2-7x+10}
\]}
{\(\left(-\infty;-2\right\rangle\cup\left\langle3;5\right)\cup\left(5;\infty\right)\)}{\(\left(-\infty;-2\right)\cup\left(3;\infty\right)\)}{\(\left(-\infty;-2\right)\cup\left(3;5\right)\cup\left(5;\infty\right)\)}{\(\left(-\infty;2\right)\cup\left(3;5\right)\cup\left(5;\infty\right)\)}{}{}}
\LANGsk{
\Question{
Určte definičný obor výrazu.
\[
\frac{\sqrt{x^2-x-6}}{x^2-7x+10}
\]}
{\(\left(-\infty;-2\right\rangle\cup\left\langle3;5\right)\cup\left(5;\infty\right)\)}{\(\left(-\infty;-2\right)\cup\left(3;\infty\right)\)}{\(\left(-\infty;-2\right)\cup\left(3;5\right)\cup\left(5;\infty\right)\)}{\(\left(-\infty;2\right)\cup\left(3;5\right)\cup\left(5;\infty\right)\)}{}{}}
\LANGes{
\Question{
Halla el dominio de la siguiente expresión.
\[
\frac{\sqrt{x^2-x-6}}{x^2-7x+10}
\]}
{\(\left(-\infty;-2\right]\cup\left[3;5\right)\cup\left(5;\infty\right)\)}{\(\left(-\infty;-2\right)\cup\left(3;\infty\right)\)}{\(\left(-\infty;-2\right)\cup\left(3;5\right)\cup\left(5;\infty\right)\)}{\(\left(-\infty;2\right)\cup\left(3;5\right)\cup\left(5;\infty\right)\)}{}{}}
\End

\Data\ID{1003020603}
\Updated{2022-08-25 10:08:13}
\Level{A}
\SubArea{Radical Equations and Inequalities}
\LANGen{
\Question{
Find the domain of the expression.
\[
\sqrt{-x^2+x+20}
\]}
{\(\left[-4;5\right]\)}{\(\left(-\infty;-4\right]\cup\left[5;\infty\right)\)}{\(\emptyset\)}{\(\left[-5;4\right]\)}{}{}}
\LANGpl{
\Question{
Znajdź dziedzinę podanego wyrażenia.
\[
\sqrt{-x^2+x+20}
\]}
{\(x\in\left\langle-4;5\right\rangle\)}{\(x\in\left(-\infty;-4\right\rangle\cup\left\langle5;\infty\right)\)}{\(x\in\emptyset\)}{\(x\in\left\langle-5;4\right\rangle\)}{}{}}
\LANGcs{
\Question{
Určete definiční obor výrazu.
\[
\sqrt{-x^2+x+20}
\]}
{\(\left\langle-4;5\right\rangle\)}{\(\left(-\infty;-4\right\rangle\cup\left\langle5;\infty\right)\)}{\(\emptyset\)}{\(\left\langle-5;4\right\rangle\)}{}{}}
\LANGsk{
\Question{
Určte definičný obor výrazu.
\[
\sqrt{-x^2+x+20}
\]}
{\(\left\langle-4;5\right\rangle\)}{\(\left(-\infty;-4\right\rangle\cup\left\langle5;\infty\right)\)}{\(\emptyset\)}{\(\left\langle-5;4\right\rangle\)}{}{}}
\LANGes{
\Question{
Halla el dominio de la siguiente expresión.
\[
\sqrt{-x^2+x+20}
\]}
{\(\left[-4;5\right]\)}{\(\left(-\infty;-4\right]\cup\left[5;\infty\right)\)}{\(\emptyset\)}{\(\left[-5;4\right]\)}{}{}}
\End

\Data\ID{1003020604}
\Updated{2022-08-25 10:08:13}
\Level{A}
\SubArea{Radical Equations and Inequalities}
\LANGen{
\Question{
Find the domain of the expression.
\[
\frac{\sqrt{-x^2-2x+24}}{2x^2-3x+3}
\]}
{\(\left[-6;4\right]\)}{\(\left[-4;6\right]\)}{\(\mathbb{R}\setminus\left[-6;4\right]\)}{\(\mathbb{R}\setminus\left[-4;6\right]\)}{}{}}
\LANGpl{
\Question{
Znajdź dziedzinę wyrażenia.
\[
\frac{\sqrt{-x^2-2x+24}}{2x^2-3x+3}
\]}
{\(x\in\left\langle-6;4\right\rangle\)}{\(x\in\left\langle-4;6\right\rangle\)}{\(x\in\mathbb{R}\setminus\left\langle-6;4\right\rangle\)}{\(x\in\mathbb{R}\setminus\left\langle-4;6\right\rangle\)}{}{}}
\LANGcs{
\Question{
Určete definiční obor výrazu.
\[
\frac{\sqrt{-x^2-2x+24}}{2x^2-3x+3}
\]}
{\(\left\langle-6;4\right\rangle\)}{\(\left\langle-4;6\right\rangle\)}{\(\mathbb{R}\setminus\left\langle-6;4\right\rangle\)}{\(\mathbb{R}\setminus\left\langle-4;6\right\rangle\)}{}{}}
\LANGsk{
\Question{
Určte definičný obor výrazu.
\[
\frac{\sqrt{-x^2-2x+24}}{2x^2-3x+3}
\]}
{\(\left\langle-6;4\right\rangle\)}{\(\left\langle-4;6\right\rangle\)}{\(\mathbb{R}\setminus\left\langle-6;4\right\rangle\)}{\(\mathbb{R}\setminus\left\langle-4;6\right\rangle\)}{}{}}
\LANGes{
\Question{
Halla el dominio de la siguiente expresión.
\[
\frac{\sqrt{-x^2-2x+24}}{2x^2-3x+3}
\]}
{\(\left[-6;4\right]\)}{\(\left[-4;6\right]\)}{\(\mathbb{R}\setminus\left[-6;4\right]\)}{\(\mathbb{R}\setminus\left[-4;6\right]\)}{}{}}
\End

\Data\ID{1003187201}
\Updated{2022-04-23 05:04:31}
\Level{A}
\SubArea{Radical Equations and Inequalities}
\LANGen{
\Question{
If \( \sqrt{x^2-8x+16}=4-x \), then the number \( x \)  can be equal to:}
{\( -2 \)}{\( 10 \)}{\( 6 \)}{\( 8 \)}{}{}}
\LANGpl{
\Question{
Jeśli  \( \sqrt{x^2-8x+16}=4-x \), wtedy liczba  \( x \)  może równać się:}
{\( -2 \)}{\( 10 \)}{\( 6 \)}{\( 8 \)}{}{}}
\LANGcs{
\Question{
Když \( \sqrt{x^2-8x+16}=4-x \), číslo \( x \) pak může být rovno:}
{\( -2 \)}{\( 10 \)}{\( 6 \)}{\( 8 \)}{}{}}
\LANGsk{
\Question{
Ak \( \sqrt{x^2-8x+16}=4-x \), potom číslo \( x \)  sa môže rovnať:}
{\( -2 \)}{\( 10 \)}{\( 6 \)}{\( 8 \)}{}{}}
\LANGes{
\Question{
Suponiendo que \( \sqrt{x^2-8x+16}=4-x \), el número \( x \)  puede ser igual a:}
{\( -2 \)}{\( 10 \)}{\( 6 \)}{\( 8 \)}{}{}}
\End

\Data\ID{2000004601}
\Updated{2023-05-30 06:05:18}
\Level{A}
\SubArea{Radical Equations and Inequalities}
\LANGen{
\Question{
Solve the following equation.
\[ \sqrt{x^2+2x+5} =x+3\]}
{\( x=-1\)}{\( x=-2\)}{\( x=2\)}{\( x=1\)}{}{}}
\LANGpl{
\Question{
Rozwiąż równanie.
\[ \sqrt{x^2+2x+5} =x+3\]}
{\( x=-1\)}{\( x=-2\)}{\( x=2\)}{\( x=1\)}{}{}}
\LANGcs{
\Question{
Řešte následující rovnici.
\[ \sqrt{x^2+2x+5} =x+3\]}
{\( x=-1\)}{\( x=-2\)}{\( x=2\)}{\( x=1\)}{}{}}
\LANGsk{
\Question{
Vyriešte následujúcu rovnicu.
\[ \sqrt{x^2+2x+5} =x+3\]}
{\( x=-1\)}{\( x=-2\)}{\( x=2\)}{\( x=1\)}{}{}}
\LANGes{
\Question{
Resuelve la siguiente ecuación. 
\[ \sqrt{x^2+2x+5} =x+3\]}
{\( x=-1\)}{\( x=-2\)}{\( x=2\)}{\( x=1\)}{}{}}
\End

\Data\ID{2000004602}
\Updated{2021-12-18 10:12:44}
\Level{A}
\SubArea{Radical Equations and Inequalities}
\LANGen{
\Question{
Choose the correct statement about the following equation.
\[ \sqrt{5x+9} = \sqrt{x+1}\]}
{The equation has no solution in \( \mathbb{R} \).}{The equation has just one root and the root is an odd number.}{The equation has just one root and the root is an even number.}{The equation has just two roots.}{}{}}
\LANGpl{
\Question{
Wybierz prawidłowe stwierdzenie dotyczące następującego równania.
\[ \sqrt{5x+9} = \sqrt{x+1}\]}
{Równanie nie ma rozwiązania w \( \mathbb{R} \).}{Równanie ma dokładnie jeden pierwiastek, a pierwiastek jest liczbą nieparzystą.}{Równanie ma dokładnie jeden pierwiastek, a pierwiastek jest liczbą parzystą.}{Równanie ma dokładnie dwa pierwiastki.}{}{}}
\LANGcs{
\Question{
Vyberte pravdivé tvrzení týkající se následující rovnice.
\[ \sqrt{5x+9} = \sqrt{x+1}\]}
{Rovnice nemá řešení v oboru \( \mathbb{R} \).}{Rovnice má jen jeden kořen, kterým je liché číslo.}{Rovnice má jen jeden kořen, kterým je sudé číslo.}{Rovnice má právě dva kořeny.}{}{}}
\LANGsk{
\Question{
Vyberte správne tvrdenie o nasledujúcej rovnici.
\[ \sqrt{5x+9} = \sqrt{x+1}\]}
{Rovnica nemá riešenie v  \( \mathbb{R} \).}{Rovnica má práve jeden koreň a koreňom je nepárne číslo.}{Rovnica má práve jeden koreň a koreňom je párne číslo.}{Rovnica má práve dva korene.}{}{}}
\LANGes{
\Question{
Elige la proposición lógica sobre la siguiente ecuación:
\[ \sqrt{5x+9} = \sqrt{x+1}\]}
{La ecuación no tiene solución en \( \mathbb{R} \).}{La ecuación tiene solo una solución y se trata de un número impar.}{La ecuación tiene solo una solución y se trata de un número par.}{La ecuación tiene solo dos soluciones.}{}{}}
\End

\Data\ID{2000004603}
\Updated{2023-05-30 06:05:48}
\Level{A}
\SubArea{Radical Equations and Inequalities}
\LANGen{
\Question{
Choose the correct statement about the following equation.
\[ 2\sqrt{x-5} =2\]}
{The equation has just one positive root.}{The equation has just one negative root.}{The equation has two roots.}{The equation has no solution in \( \mathbb{R}\).}{}{}}
\LANGpl{
\Question{
Wybierz prawidłowe stwierdzenie dotyczące następującego równania.
\[ 2\sqrt{x-5} =2\]}
{Równanie ma dokładnie jeden dodatni pierwiastek.}{Równanie ma dokładnie jeden ujemny pierwiastek.}{Równanie ma dwa pierwiastki.}{Równanie nie ma rozwiązania w \( \mathbb{R}\).}{}{}}
\LANGcs{
\Question{
Vyberte pravdivé tvrzení o následující rovnici.
\[ 2\sqrt{x-5} =2\]}
{Rovnice má právě jeden kladný kořen.}{Rovnice má právě jeden záporný kořen.}{Rovnice má dva kořeny.}{Rovnice nemá v \( \mathbb{R}\) řešení.}{}{}}
\LANGsk{
\Question{
Vyberte správne tvrdenie o nasledujúcej rovnici.
\[ 2\sqrt{x-5} =2\]}
{Rovnica má práve jeden kladný koreň.}{Rovnica má práve jeden záporný koreň.}{Rovnica má dva korene.}{Rovnica nemá riešenie v \( \mathbb{R}\).}{}{}}
\LANGes{
\Question{
Elige la proposición lógica verdadera de la siguiente ecuación. 
\[ 2\sqrt{x-5} =2\]}
{La ecuación tiene una raíz positiva.}{La ecuación tiene una raíz negativa.}{La ecuación tiene dos raíces.}{La ecuación no tiene solución en \( \mathbb{R}\).}{}{}}
\End

\Data\ID{2000004604}
\Updated{2023-05-30 06:05:01}
\Level{A}
\SubArea{Radical Equations and Inequalities}
\LANGen{
\Question{
Choose the correct statement about the following equation.
\[ 4+ 2\sqrt{x+4} =x\]}
{The root of the equation is \(x=12\).}{The equation has no root in \( \mathbb{R}\).}{The roots of the equation are \(x_1=0\) and \(x_2=12\).}{The roots of the equation are \(x_1=4\) and \(x_2= -2\).}{}{}}
\LANGpl{
\Question{
Wybierz prawidłowe stwierdzenie dotyczące następującego równania.
\[ 4+ 2\sqrt{x+4} =x\]}
{Pierwiastkiem równania jest \(x=12\).}{Równanie nie ma pierwiastka w \( \mathbb{R}\).}{Pierwiastkami równania są \(x_1=0\) i \(x_2=12\).}{Pierwiastkami równania są \(x_1=4\) i \(x_2= -2\).}{}{}}
\LANGcs{
\Question{
Vyberte pravdivé tvrzení o následující rovnici.
\[ 4+ 2\sqrt{x+4} =x\]}
{Kořenem rovnice je \(x=12\).}{Rovnice nemá v \( \mathbb{R}\) řešení.}{Kořeny rovnice jsou \(x_1=0\) a \(x_2=12\).}{Kořeny rovnice jsou \(x_1=4\) a \(x_2= -2\).}{}{}}
\LANGsk{
\Question{
Vyberte správne tvrdenie o nasledujúcej rovnici.
\[ 4+ 2\sqrt{x+4} =x\]}
{Koreňom  rovnice je  \(x=12\).}{Rovnica nemá riešenie v  \( \mathbb{R}\).}{Koreňmi rovnice sú  \(x_1=0\) a \(x_2=12\).}{Koreňmi rovnice sú  \(x_1=4\) a \(x_2= -2\).}{}{}}
\LANGes{
\Question{
Elige la proposición lógica verdadera de la siguiente ecuación. 
\[ 4+ 2\sqrt{x+4} =x\]}
{La raíz de la ecuación es \(x=12\).}{La ecuación no tiene ninguna raíz en \( \mathbb{R}\).}{Las raíces de la ecuación son \(x_1=0\) y \(x_2=12\).}{Las raíces de la ecuación son \(x_1=4\) y \(x_2= -2\).}{}{}}
\End

\Data\ID{2000004605}
\Updated{2023-05-30 06:05:28}
\Level{A}
\SubArea{Radical Equations and Inequalities}
\LANGen{
\Question{
Choose the correct statement about the following equation.
\[ \sqrt{4-x} = 3-\sqrt{x-4}\]}
{The equation has no solution in \( \mathbb{R}\).}{The equation has just one root and the root is an odd number.}{The equation has just one root and the root is an even number.}{The equation has just two roots.}{}{}}
\LANGpl{
\Question{
Wybierz prawidłowe stwierdzenie dotyczące następującego równania.
\[ \sqrt{4-x} = 3-\sqrt{x-4}\]}
{Równanie nie ma rozwiązania w  \( \mathbb{R}\).}{Równanie ma tylko jeden pierwiastek, a pierwiastek jest liczbą nieparzystą.}{Równanie ma tylko jeden pierwiastek, a pierwiastek jest liczbą parzystą.}{Równanie ma dokładnie dwa pierwiastki.}{}{}}
\LANGcs{
\Question{
Vyberte pravdivé tvrzení o následující rovnici.
\[ \sqrt{4-x} = 3-\sqrt{x-4}\]}
{Rovnice nemá v \( \mathbb{R}\) řešení.}{Rovnice má v \( \mathbb{R}\) právě jeden kořen a je jím liché číslo.}{Rovnice má v \( \mathbb{R}\) právě jeden kořen a je jím sudé číslo.}{Rovnice má právě dva reálné kořeny.}{}{}}
\LANGsk{
\Question{
Vyberte správne tvrdenie o nasledujúcej rovnici.
\[ \sqrt{4-x} = 3-\sqrt{x-4}\]}
{Rovnica nemá riešenie v \( \mathbb{R}\).}{Rovnica má iba jeden koreň a koreňom je nepárne číslo.}{Rovnica má iba jeden koreň a koreňom  je párne číslo.}{Rovnica má iba dva korene.}{}{}}
\LANGes{
\Question{
Elige la proposición lógica verdadera de la siguiente ecuación. 
\[ \sqrt{4-x} = 3-\sqrt{x-4}\]}
{La ecuación no tiene solución en \( \mathbb{R}\).}{La ecuación tiene una raíz y esta raíz es un número impar.}{La ecuación tiene una raíz y la raíz es un número par.}{La ecuación tiene dos raíces.}{}{}}
\End

\Data\ID{2000004606}
\Updated{2023-05-30 06:05:43}
\Level{A}
\SubArea{Radical Equations and Inequalities}
\LANGen{
\Question{
Choose the correct statement about the following equation.
\[ 5+ 2\sqrt{x-2}=9\]}
{The solution of this equation is a number from  the interval \( (5;6] \).}{The solution of this equation is a number from  the interval \( (-1 ; 0 ] \).}{The solution of this equation is a number from  the interval \( (0;3] \).}{The solution of this equation is a number from  the interval \( (3;5 ] \).}{}{}}
\LANGpl{
\Question{
Wybierz prawidłowe stwierdzenie dotyczące następującego równania.
\[ 5+ 2\sqrt{x-2}=9\]}
{Rozwiązaniem tego równania jest liczba należąca do przedziału \( (5;6\rangle \).}{Rozwiązaniem tego równania jest liczba należąca do przedziału \( (-1 ; 0 \rangle \).}{Rozwiązaniem tego równania jest liczba należąca do przedziału \( (0;3\rangle \).}{Rozwiązaniem tego równania jest liczba należąca do przedziału \( (3;5 \rangle \).}{}{}}
\LANGcs{
\Question{
Vyberte pravdivé tvrzení o následující rovnici. 
\[ 5+ 2\sqrt{x-2}=9\]}
{Řešením rovnice je číslo z intervalu \( (5;6\rangle \).}{Řešením rovnice je číslo z intervalu \( (-1 ; 0 \rangle \).}{Řešením rovnice je číslo z intervalu \( (0;3\rangle \).}{Řešením rovnice je číslo z intervalu \( (3;5 \rangle \).}{}{}}
\LANGsk{
\Question{
Vyberte správne tvrdenie o nasledujúcej rovnici.
\[ 5+ 2\sqrt{x-2}=9\]}
{Riešením tejto rovnice je číslo patriace do intervalu  \( (5;6\rangle \).}{Riešením tejto rovnice je číslo patriace do intervalu  \( (-1 ; 0 \rangle \).}{Riešením tejto rovnice je číslo patriace do intervalu  \( (0;3\rangle \).}{Riešením tejto rovnice je číslo patriace do intervalu \( (3;5 \rangle \).}{}{}}
\LANGes{
\Question{
Elige la proposición lógica verdadera de la siguiente ecuación. 
\[ 5+ 2\sqrt{x-2}=9\]}
{La solución de la ecuación es un número del intervalo \( (5;6] \).}{La solución de la ecuación es un número del intervalo \( (-1 ; 0 ] \).}{La solución de la ecuación es un número del intervalo \( (0;3] \).}{La solución de la ecuación es un número del intervalo \( (3;5 ] \).}{}{}}
\End

\Data\ID{2000004607}
\Updated{2023-05-30 06:05:57}
\Level{A}
\SubArea{Radical Equations and Inequalities}
\LANGen{
\Question{
Choose the correct statement about the following equation.
\[ x-2 \sqrt{x-6} =6 \]}
{The roots of the equation are \(x_1=10\) and \(x_2= 6\).}{The only root of the equation is \(x=6\).}{The equation has no root in \( \mathbb{R}\).}{The roots of the equation are \(x_1=-2\) and \(x_2=6\).}{}{}}
\LANGpl{
\Question{
Wybierz prawidłowe stwierdzenie dotyczące następującego równania.
\[ x-2 \sqrt{x-6} =6 \]}
{Pierwiastkami równania są \(x_1=10\) i \(x_2= 6\).}{Jedynym pierwiastkiem równania jest \(x=6\).}{Równanie nie ma pierwiastka w \( \mathbb{R}\).}{Pierwiastkami równania są \(x_1=-2\) i \(x_2=6\).}{}{}}
\LANGcs{
\Question{
Vyberte pravdivé tvrzení o následující rovnici. 
\[ x-2 \sqrt{x-6} =6 \]}
{Kořeny rovnice jsou  \(x_1=10\) a \(x_2= 6\).}{Jediným kořenem rovnice je \(x=6\).}{Rovnice nemá v \( \mathbb{R}\) žádný kořen.}{Kořeny rovnice jsou \(x_1=-2\) a \(x_2=6\).}{}{}}
\LANGsk{
\Question{
Vyberte správne tvrdenie o nasledujúcej rovnici.
\[ x-2 \sqrt{x-6} =6 \]}
{Korene rovnice sú \(x_1=10\) a \(x_2= 6\).}{Jediným koreňom rovnice je \(x=6\).}{Rovnica nemá koreň v \( \mathbb{R}\).}{Korene rovnice sú  \(x_1=-2\)  a  \(x_2=6\).}{}{}}
\LANGes{
\Question{
Elige la proposición lógica verdadera de la siguiente ecuación. 
\[ x-2 \sqrt{x-6} =6 \]}
{Las raíces de la ecuación son \(x_1=10\) y \(x_2= 6\).}{La única raíz de la ecuación es \(x=6\).}{La ecuación no tiene raíces en \( \mathbb{R}\).}{Las raíces de la ecuación son \(x_1=-2\) y \(x_2=6\).}{}{}}
\End

\Data\ID{2000004608}
\Updated{2023-05-30 06:05:14}
\Level{A}
\SubArea{Radical Equations and Inequalities}
\LANGen{
\Question{
Choose the correct statement about the following equation.
\[
\sqrt{x+5} + \sqrt{2-x}=0
\]}
{The equation has no solution in \(\mathbb{R}\).}{The equation has just one root and the root is an odd number.}{The equation has just one root and the root is an even number.}{The equation has just two roots.}{}{}}
\LANGpl{
\Question{
Wybierz prawidłowe stwierdzenie dotyczące następującego równania.
\[
\sqrt{x+5} + \sqrt{2-x}=0
\]}
{Równanie nie ma rozwiązania w \(\mathbb{R}\).}{Równanie ma dokładnie jeden pierwiastek, a pierwiastek jest liczbą nieparzystą.}{Równanie ma dokładnie jeden pierwiastek, a pierwiastek jest liczbą parzystą.}{Równanie ma dokładnie dwa pierwiastki.}{}{}}
\LANGcs{
\Question{
Vyberte pravdivé tvrzení o následující rovnici. 
\[
\sqrt{x+5} + \sqrt{2-x}=0
\]}
{Rovnice nemá v \(\mathbb{R}\) žádné řešení.}{Rovnice má právě jedno řešení a je jím liché číslo.}{Rovnice má právě jedno řešení a je jím sudé číslo.}{Rovnice má právě dva různé kořeny.}{}{}}
\LANGsk{
\Question{
Vyberte správne tvrdenie o nasledujúcej rovnici.
\[
\sqrt{x+5} + \sqrt{2-x}=0
\]}
{Rovnica nemá riešenie v  \(\mathbb{R}\).}{Rovnica má práve jeden koreň a koreňom je nepárne číslo.}{Rovnica má práve jeden koreň a koreňom je párne číslo.}{Rovnica má práve dva korene.}{}{}}
\LANGes{
\Question{
Elige la proposición lógica verdadera de la siguiente ecuación. 
\[\sqrt{x+5} + \sqrt{2-x}=0
\]}
{La ecuación no tiene solución en \(\mathbb{R}\).}{La ecuación tiene una raíz y esta raíz es un número impar.}{La ecuación tiene una raíz y esta raíz es un número par.}{La ecuación tiene dos raíces.}{}{}}
\End

\Data\ID{2000004609}
\Updated{2023-05-30 06:05:28}
\Level{A}
\SubArea{Radical Equations and Inequalities}
\LANGen{
\Question{
Choose the correct statement about the following equation.
\[
\sqrt{7-\sqrt{x-3}}=2
\]}
{The equation has just one root and the root is an even number.}{The equation has no solution in \( \mathbb{R}\).}{The equation has just one root and the root is an odd number.}{The equation has just two roots.}{}{}}
\LANGpl{
\Question{
Wybierz prawidłowe stwierdzenie dotyczące następującego równania.
\[
\sqrt{7-\sqrt{x-3}}=2
\]}
{Równanie ma dokładnie jeden pierwiastek, a pierwiastek jest liczbą parzystą.}{Równanie nie ma rozwiązania w \( \mathbb{R}\).}{Równanie ma dokładnie jeden pierwiastek, a pierwiastek jest liczbą nieparzystą.}{Równanie ma dokładnie dwa pierwiastki.}{}{}}
\LANGcs{
\Question{
Vyberte pravdivé tvrzení o následující rovnici. 
\[
\sqrt{7-\sqrt{x-3}}=2
\]}
{Rovnice má právě jedno řešení a je jím sudé číslo.}{Rovnice nemá v \(\mathbb{R}\) žádné řešení.}{Rovnice má právě jedno řešení a je jím liché číslo.}{Rovnice má právě dva různé kořeny.}{}{}}
\LANGsk{
\Question{
Vyberte správne tvrdenie o nasledujúcej rovnici.
\[
\sqrt{7-\sqrt{x-3}}=2
\]}
{Rovnica má iba jeden koreň a koreň je párne číslo.}{Rovnica nemá riešenie v  \( \mathbb{R}\).}{Rovnica má iba jeden koreň a koreň je nepárne číslo.}{Rovnica má iba dva korene.}{}{}}
\LANGes{
\Question{
Elige la proposición lógica verdadera de la siguiente ecuación. 
\[
\sqrt{7-\sqrt{x-3}}=2
\]}
{La ecuación tiene una raíz y esta raíz es un número par.}{La ecuación no tiene solución en \( \mathbb{R}\).}{La ecuación tiene una raíz y esta raíz es un número impar.}{La ecuación tiene dos raíces.}{}{}}
\End

\Data\ID{2000004610}
\Updated{2023-05-30 06:05:41}
\Level{A}
\SubArea{Radical Equations and Inequalities}
\LANGen{
\Question{
Choose the correct statement about the following equation.
\[ \sqrt{(x+5)^2} =x+5\]}
{The solution of the equation are all \( x \in [ -5; \infty) \).}{The solution of the equation are all \( x \in \mathbb{R} \).}{The equation has no solution in \(\mathbb{R} \).}{The equation has just one root \( x=0\).}{}{}}
\LANGpl{
\Question{
Wskaż prawidłowe stwierdzenie dotyczące następującego równania.
\[ \sqrt{(x+5)^2} =x+5\]}
{Rozwiązaniem równania są wszystkie \( x \in \langle -5; \infty) \).}{Rozwiązaniem równania są wszystkie \( x \in \mathbb{R} \).}{Równanie nie ma rozwiązania w \(\mathbb{R} \).}{Równanie ma dokładnie jeden pierwiastek  \( x=0\).}{}{}}
\LANGcs{
\Question{
Vyberte pravdivé tvrzení o následující rovnici. 
\[ \sqrt{(x+5)^2} =x+5\]}
{Řešením rovnice jsou všechna \( x \in \langle -5; \infty) \).}{Řešením rovnice jsou všechna \( x \in \mathbb{R} \).}{Rovnice nemá v \(\mathbb{R}\) žádné řešení.}{Rovnice má právě jeden kořen \( x=0\).}{}{}}
\LANGsk{
\Question{
Vyberte správne tvrdenie o nasledujúcej rovnici.
\[ \sqrt{(x+5)^2} =x+5\]}
{Riešením rovnice sú všetky \( x \in \langle -5; \infty) \).}{Riešením sú všetky  \( x \in \mathbb{R} \).}{Rovnica nemá riešenie v  \(\mathbb{R} \).}{Rovnica má práve jeden koreň  \( x=0\).}{}{}}
\LANGes{
\Question{
Elige la proposición lógica verdadera sobre la siguiente ecuación. 
\[ \sqrt{(x+5)^2} =x+5\]}
{La solución de la ecuación son todos los \( x \in [ -5; \infty) \).}{La solución de la ecuación son todos los \( x \in \mathbb{R} \).}{La ecuación no tiene solución en \(\mathbb{R} \).}{La ecuación tiene la única raíz \( x=0\).}{}{}}
\End

\Data\ID{2010007701}
\Updated{2022-02-15 07:02:37}
\Level{A}
\SubArea{Radical Equations and Inequalities}
\LANGen{
\Question{
Find the solution of the equation.
\[
\sqrt{-2x^2+3x-4}=2
\]}
{\(x\in\emptyset\)}{\(x\in\mathbb{R}\)}{\(x\in (-3;3)\)}{\(x\in\{-3;3\}\)}{}{}}
\LANGpl{
\Question{
Rozwiąż równanie.
\[
\sqrt{-2x^2+3x-4}=2
\]}
{\(x\in\emptyset\)}{\(x\in\mathbb{R}\)}{\(x\in (-3;3)\)}{\(x\in\{-3;3\}\)}{}{}}
\LANGcs{
\Question{
Najděte řešení rovnice.
\[
\sqrt{-2x^2+3x-4}=2
\]}
{\(x\in\emptyset\)}{\(x\in\mathbb{R}\)}{\(x\in (-3;3)\)}{\(x\in\{-3;3\}\)}{}{}}
\LANGsk{
\Question{
Nájdite riešenie rovnice.
\[
\sqrt{-2x^2+3x-4}=2
\]}
{\(x\in\emptyset\)}{\(x\in\mathbb{R}\)}{\(x\in (-3;3)\)}{\(x\in\{-3;3\}\)}{}{}}
\LANGes{
\Question{
Halla la solución de la siguiente ecuación.
\[
\sqrt{-2x^2+3x-4}=2
\]}
{\(x\in\emptyset\)}{\(x\in\mathbb{R}\)}{\(x\in (-3;3)\)}{\(x\in\{-3;3\}\)}{}{}}
\End

\Data\ID{2010007702}
\Updated{2022-11-01 09:11:36}
\Level{A}
\SubArea{Radical Equations and Inequalities}
\LANGen{
\Question{
Find the solution of the equation.
\[
\sqrt{x^2-3x+6}=\sqrt{x^2-4x+6}
\]}
{\(x\in\{0\}\)}{\(x\in\emptyset\)}{\(x\in\mathbb{R}\)}{\(x\in\mathbb{R}\setminus\{0\}\)}{}{}}
\LANGpl{
\Question{
Rozwiąż równanie.
\[
\sqrt{x^2-3x+6}=\sqrt{x^2-4x+6}
\]}
{\(x\in\{0\}\)}{\(x\in\emptyset\)}{\(x\in\mathbb{R}\)}{\(x\in\mathbb{R}\setminus\{0\}\)}{}{}}
\LANGcs{
\Question{
Co je řešením rovnice?
\[
\sqrt{x^2-3x+6}=\sqrt{x^2-4x+6}
\]}
{\(x\in\{0\}\)}{\(x\in\emptyset\)}{\(x\in\mathbb{R}\)}{\(x\in\mathbb{R}\setminus\{0\}\)}{}{}}
\LANGsk{
\Question{
Nájdite riešenie rovnice.
\[
\sqrt{x^2-3x+6}=\sqrt{x^2-4x+6}
\]}
{\(x\in\{0\}\)}{\(x\in\emptyset\)}{\(x\in\mathbb{R}\)}{\(x\in\mathbb{R}\setminus\{0\}\)}{}{}}
\LANGes{
\Question{
Halla la solución de la siguiente ecuación.
\[
\sqrt{x^2-3x+6}=\sqrt{x^2-4x+6}
\]}
{\(x\in\{0\}\)}{\(x\in\emptyset\)}{\(x\in\mathbb{R}\)}{\(x\in\mathbb{R}\setminus\{0\}\)}{}{}}
\End

\Data\ID{2010007703}
\Updated{2022-02-15 07:02:37}
\Level{A}
\SubArea{Radical Equations and Inequalities}
\LANGen{
\Question{
Find the domain of the expression.
\[
\frac1{\sqrt{2x^2-11x+14}}
\]}
{\(\left(-\infty;2\right)\cup\left(\frac72;\infty\right)\)}{\(\left(2;\frac72\right)\)}{\(\left(-\infty;2\right] \cup \left[ \frac72;\infty\right)\)}{\(\left[ 2;\frac72 \right] \)}{}{}}
\LANGpl{
\Question{
Znajdź dziedzinę wyrażenia.
\[
\frac1{\sqrt{2x^2-11x+14}}
\]}
{\(\left(-\infty;2\right)\cup\left(\frac72;\infty\right)\)}{\(\left(2;\frac72\right)\)}{\(\left(-\infty;2\right\rangle \cup \left\langle \frac72;\infty\right)\)}{\(\left\langle 2;\frac72 \right\rangle \)}{}{}}
\LANGcs{
\Question{
Vyberte definiční obor výrazu.
\[
\frac1{\sqrt{2x^2-11x+14}}
\]}
{\(\left(-\infty;2\right)\cup\left(\frac72;\infty\right)\)}{\(\left(2;\frac72\right)\)}{\(\left(-\infty;2\right\rangle \cup \left\langle \frac72;\infty\right)\)}{\(\left\langle 2;\frac72 \right\rangle \)}{}{}}
\LANGsk{
\Question{
Nájdite definičný obor výrazu.
\[
\frac1{\sqrt{2x^2-11x+14}}
\]}
{\(\left(-\infty;2\right)\cup\left(\frac72;\infty\right)\)}{\(\left(2;\frac72\right)\)}{\(\left(-\infty;2\right\rangle \cup \left\langle \frac72;\infty\right)\)}{\(\left\langle 2;\frac72 \right\rangle \)}{}{}}
\LANGes{
\Question{
Halla el dominio de la siguiente expresión.
\[
\frac1{\sqrt{2x^2-11x+14}}
\]}
{\(\left(-\infty;2\right)\cup\left(\frac72;\infty\right)\)}{\(\left(2;\frac72\right)\)}{\(\left(-\infty;2\right] \cup \left[ \frac72;\infty\right)\)}{\(\left[ 2;\frac72 \right] \)}{}{}}
\End

\Data\ID{2010007704}
\Updated{2022-11-01 09:11:20}
\Level{A}
\SubArea{Radical Equations and Inequalities}
\LANGen{
\Question{
Find the domain of the expression.
\[
\sqrt{-x^2-7x+30}
\]}
{\([ -10;3] \)}{\(\left(-\infty;-10 ] \cup[ 3;\infty\right)\)}{\([ -3;10] \)}{\(\emptyset\)}{}{}}
\LANGpl{
\Question{
Znajdź dziedzinę wyrażenia.
\[
\sqrt{-x^2-7x+30}
\]}
{\(\langle -10;3\rangle \)}{\(\left(-\infty;-10 \rangle \cup\langle 3;\infty\right)\)}{\(\langle -3;10\rangle \)}{\(\emptyset\)}{}{}}
\LANGcs{
\Question{
Najděte definiční obor výrazu.
\[
\sqrt{-x^2-7x+30}
\]}
{\(\langle -10;3\rangle \)}{\(\left(-\infty;-10 \rangle \cup\langle 3;\infty\right)\)}{\(\langle -3;10\rangle \)}{\(\emptyset\)}{}{}}
\LANGsk{
\Question{
Nájdite definičný obor výrazu.
\[
\sqrt{-x^2-7x+30}
\]}
{\(\langle -10;3\rangle \)}{\(\left(-\infty;-10 \rangle \cup\langle 3;\infty\right)\)}{\(\langle -3;10\rangle \)}{\(\emptyset\)}{}{}}
\LANGes{
\Question{
Halla el dominio de la siguiente expresión.
\[
\sqrt{-x^2-7x+30}
\]}
{\([ -10;3] \)}{\(\left(-\infty;-10 ] \cup[ 3;\infty\right)\)}{\([ -3;10] \)}{\(\emptyset\)}{}{}}
\End

\Data\ID{2010007706}
\Updated{2022-02-15 07:02:37}
\Level{A}
\SubArea{Radical Equations and Inequalities}
\LANGen{
\Question{
If \( \sqrt{x^2-x-12}=x+3 \), then the number \( x \)  can be equal to:}
{\( -3\)}{\( -4\)}{\( -5\)}{\( -6\)}{}{}}
\LANGpl{
\Question{
Jeżeli \( \sqrt{x^2-x-12}=x+3 \), to liczba \( x \)  jest równa:}
{\( -3\)}{\( -4\)}{\( -5\)}{\( -6\)}{}{}}
\LANGcs{
\Question{
Čemu může být rovno \( x \), pokud platí, že \( \sqrt{x^2-x-12}=x+3 \)?}
{\( -3\)}{\( -4\)}{\( -5\)}{\( -6\)}{}{}}
\LANGsk{
\Question{
Ak platí \( \sqrt{x^2-x-12}=x+3 \), potom číslo \( x \) sa môže rovnať:}
{\( -3\)}{\( -4\)}{\( -5\)}{\( -6\)}{}{}}
\LANGes{
\Question{
Suponiendo que \( \sqrt{x^2-x-12}=x+3 \), el número \( x \)  puede ser igual a:}
{\( -3\)}{\( -4\)}{\( -5\)}{\( -6\)}{}{}}
\End

\Data\ID{2010007707}
\Updated{2022-11-01 09:11:05}
\Level{A}
\SubArea{Radical Equations and Inequalities}
\LANGen{
\Question{
Find the domain of the following equation. 
\[ 
 \sqrt{7 - x} = 13
\]}
{\((-\infty ;7 ]  \)}{\( [ -7;\infty )\)}{\((-\infty ;6 )  \)}{\( (6;\infty )\)}{}{}}
\LANGpl{
\Question{
Znajdź dziedzinę równania. 
\[ 
 \sqrt{7 - x} = 13
\]}
{\((-\infty ;7 \rangle  \)}{\( \langle -7;\infty )\)}{\((-\infty ;6 )  \)}{\( (6;\infty )\)}{}{}}
\LANGcs{
\Question{
Jaký je definiční obor následující rovnice? 
\[ 
 \sqrt{7 - x} = 13
\]}
{\((-\infty ;7 \rangle  \)}{\( \langle -7;\infty )\)}{\((-\infty ;6 )  \)}{\( (6;\infty )\)}{}{}}
\LANGsk{
\Question{
Nájdite definičný obor rovnice.
\[ 
 \sqrt{7 - x} = 13
\]}
{\((-\infty ;7 \rangle  \)}{\( \langle -7;\infty )\)}{\((-\infty ;6 )  \)}{\( (6;\infty )\)}{}{}}
\LANGes{
\Question{
Halla el dominio de la siguiente ecuación.
\[ 
 \sqrt{7 - x} = 13
\]}
{\((-\infty ;7 ]  \)}{\( [ -7;\infty )\)}{\((-\infty ;6 )  \)}{\( (6;\infty )\)}{}{}}
\End

\Data\ID{2010007708}
\Updated{2022-02-15 07:02:34}
\Level{A}
\SubArea{Radical Equations and Inequalities}
\LANGen{
\Question{
Find the domain of the following equation. 
\[ 
 \sqrt{9-3x} + \sqrt{2x+8} = 13
\]}
{\( [ -4;3 ] \)}{\((-\infty ;-4] \)}{\( [ -4;\infty )\)}{\( [ 3;\infty )\)}{}{}}
\LANGpl{
\Question{
Znajdź dziedzinę równania. 
\[ 
 \sqrt{9-3x} + \sqrt{2x+8} = 13
\]}
{\( \langle -4;3 \rangle \)}{\((-\infty ;-4\rangle \)}{\( \langle -4;\infty )\)}{\( \langle 3;\infty )\)}{}{}}
\LANGcs{
\Question{
Jaký je definiční obor následující rovnice? 
\[ 
 \sqrt{9-3x} + \sqrt{2x+8} = 13
\]}
{\( \langle -4;3 \rangle \)}{\((-\infty ;-4\rangle \)}{\( \langle -4;\infty )\)}{\( \langle 3;\infty )\)}{}{}}
\LANGsk{
\Question{
Nájdite definičný obor rovnice.
\[ 
 \sqrt{9-3x} + \sqrt{2x+8} = 13
\]}
{\( \langle -4;3 \rangle \)}{\((-\infty ;-4\rangle \)}{\( \langle -4;\infty )\)}{\( \langle 3;\infty )\)}{}{}}
\LANGes{
\Question{
Halla el dominio de la siguiente ecuación.
\[ 
 \sqrt{9-3x} + \sqrt{2x+8} = 13
\]}
{\( [ -4;3 ] \)}{\((-\infty ;-4] \)}{\( [ -4;\infty )\)}{\( [ 3;\infty )\)}{}{}}
\End

\Data\ID{2010007709}
\Updated{2022-11-01 09:11:52}
\Level{A}
\SubArea{Radical Equations and Inequalities}
\LANGen{
\Question{
Identify a true statement referring to the solution of the following equation.
\[ 
 \sqrt{3x - 6} = 3
\]}
{The solution is an odd number.}{The solution is a number divisible by the number
\(3\).}{The solution is an irrational number.}{The solution is not a prime number.}{}{}}
\LANGpl{
\Question{
Wskaż prawdziwe stwierdzenie dotyczące rozwiązania poniższego równania.
\[ 
 \sqrt{3x - 6} = 3
\]}
{Rozwiązaniem jest liczba nieparzysta.}{Rozwiązaniem jest liczba podzielna przez liczbę
\(3\).}{Rozwiązaniem jest liczba niewymierna.}{Rozwiązaniem nie jest liczba pierwsza.}{}{}}
\LANGcs{
\Question{
Které tvrzení týkající se řešení následující rovnice je pravdivé?
\[ 
 \sqrt{3x - 6} = 3
\]}
{Řešením je liché číslo.}{Řešením je číslo dělitelné číslem
\(3\).}{Řešením je iracionální číslo.}{Řešením není prvočíslo.}{}{}}
\LANGsk{
\Question{
Vyberte pravdivé tvrdenie týkajúce sa riešenia nasledujúcej rovnice.
\[ 
 \sqrt{3x - 6} = 3
\]}
{Koreňom rovnice je nepárne číslo.}{Koreň rovnice je deliteľný číslom 
\(3\).}{Koreňom rovnice je iracionálne číslo.}{Koreň rovnice nie je prvočíslo.}{}{}}
\LANGes{
\Question{
Identifica la proposición lógica que se refiera a la solución de la siguiente ecuación:
\[ 
 \sqrt{3x - 6} = 3
\]}
{La solución es un número impar.}{La solución es un número divisible por
\(3\).}{La solución es un número irracional.}{La solución no es un número primo.}{}{}}
\End

\Data\ID{2010007710}
\Updated{2022-11-01 09:11:40}
\Level{A}
\SubArea{Radical Equations and Inequalities}
\LANGen{
\Question{
Identify a true statement which concerns the following pair of equations.
\[ 
\begin{aligned}
 \sqrt{x+3} & = 5 &\text{(1)}
 \\
\sqrt{11-x} & = 3 &\text{(2)}
 \end{aligned}
\]}
{The solution of (1) is bigger than the solution of (2).}{The solution of (1) is smaller than the solution of (2).}{The solutions of both equations are prime numbers.}{The solution of (1) equals to the solution of (2).}{}{}}
\LANGpl{
\Question{
Wskaż prawdziwe stwierdzenie, które dotyczy następującej pary równań.
\[ 
\begin{aligned}
 \sqrt{x+3} & = 5 &\text{(1)}
 \\
\sqrt{11-x} & = 3 &\text{(2)}
 \end{aligned}
\]}
{Rozwiązanie (1) jest większe niż rozwiązanie (2).}{Rozwiązanie (1) jest mniejsze niż rozwiązanie (2).}{Rozwiązania obu równań są liczbami pierwszymi.}{Rozwiązanie (1) równa się rozwiązaniu (2).}{}{}}
\LANGcs{
\Question{
Které tvrzení týkající se následujících rovnic je pravdivé?
\[ 
\begin{aligned}
 \sqrt{x+3} & = 5 &\text{(1)}
 \\
\sqrt{11-x} & = 3 &\text{(2)}
 \end{aligned}
\]}
{Řešením rovnice (1) je větší číslo než řešením rovnice (2).}{Řešením rovnice (1) je menší číslo než řešením rovnice (2).}{Řešením obou rovnic jsou prvočísla.}{Řešení rovnice (1) je rovno řešení rovnice (2).}{}{}}
\LANGsk{
\Question{
Vyberte pravdivé tvrdenie, ktoré sa týka nasledujúcej dvojice rovníc.
\[ 
\begin{aligned}
 \sqrt{x+3} & = 5 &\text{(1)}
 \\
\sqrt{11-x} & = 3 &\text{(2)}
 \end{aligned}
\]}
{Koreň rovnice  (1) je väčší ako koreň rovnice  (2).}{Koreň rovnice  (1) je menší ako koreň rovnice  (2).}{Korene obidvoch rovníc sú prvočísla.}{Koreň rovnice  (1) sa rovná koreňu rovnice  (2).}{}{}}
\LANGes{
\Question{
Dadas las ecuaciones:
\[ 
\begin{aligned}
 \sqrt{x+3} & = 5 &\text{(1)}
 \\
\sqrt{11-x} & = 3 &\text{(2)}
 \end{aligned}
\]
Identifica la proposición lógica que hace referencia a las ecuaciones.}
{La solución de la ecuación (1) es mayor que la solución de la ecuación (2).}{La solución de la ecuación (1) es menor que la solución de la ecuación (2).}{Las soluciones de ambas ecuaciones son números primos.}{La solución de la ecuación (1) equivale a la solución de la ecuación (2).}{}{}}
\End

\Data\ID{2010007801}
\Updated{2022-11-01 09:11:06}
\Level{A}
\SubArea{Radical Equations and Inequalities}
\LANGen{
\Question{
Identify a true statement which concerns to the following equation.
\[ 
 2\sqrt{x+5} = x+2
\]}
{The solution is in the set \(\left \{x\in \mathbb{R} : -1 < x\leq - 5 \right \}\).}{The solution is in the set \(\left \{x\in \mathbb{R} : 5 < x\leq 7 \right \}\).}{The solution is in the set \(\left \{x\in \mathbb{R} : -4 < x\leq - 1 \right \}\).}{The solution is in the set \(\left \{x\in \mathbb{R} : -1 < x\leq 2 \right \}\).}{}{}}
\LANGpl{
\Question{
Wskaż prawdziwe stwierdzenie, które dotyczy poniższego równania.
\[ 
 2\sqrt{x+5} = x+2
\]}
{Rozwiązanie należy do zbioru \(\left \{x\in \mathbb{R} : -1 < x\leq - 5 \right \}\).}{Rozwiązanie należy do zbioru \(\left \{x\in \mathbb{R} : 5 < x\leq 7 \right \}\).}{Rozwiązanie należy do zbioru \(\left \{x\in \mathbb{R} : -4 < x\leq - 1 \right \}\).}{Rozwiązanie należy do zbioru \(\left \{x\in \mathbb{R} : -1 < x\leq 2 \right \}\).}{}{}}
\LANGcs{
\Question{
Které z následujících tvrzení o dané rovnici je pravdivé?
\[ 
 2\sqrt{x+5} = x+2
\]}
{Řešení se nachází v množině \(\left \{x\in \mathbb{R} : -1 < x\leq - 5 \right \}\).}{Řešení se nachází v množině \(\left \{x\in \mathbb{R} : 5 < x\leq 7 \right \}\).}{Řešení se nachází v množině \(\left \{x\in \mathbb{R} : -4 < x\leq - 1 \right \}\).}{Řešení se nachází v množině \(\left \{x\in \mathbb{R} : -1 < x\leq 2 \right \}\).}{}{}}
\LANGsk{
\Question{
Nájdite pravdivé tvrdenie, ktoré sa týka nasledujúcej rovnice.
\[ 
 2\sqrt{x+5} = x+2
\]}
{Riešením tejto rovnice je číslo z množiny  \(\left \{x\in \mathbb{R} : -1 < x\leq - 5 \right \}\).}{Riešením tejto rovnice je číslo z množiny  \(\left \{x\in \mathbb{R} : 5 < x\leq 7 \right \}\).}{Riešením tejto rovnice je číslo z množiny  \(\left \{x\in \mathbb{R} : -4 < x\leq - 1 \right \}\).}{Riešením tejto rovnice je číslo z množiny \(\left \{x\in \mathbb{R} : -1 < x\leq 2 \right \}\).}{}{}}
\LANGes{
\Question{
Dada la ecuación:
\[ 
 2\sqrt{x+5} = x+2
\]
Identifica la proposición lógica que hace referencia a la ecuación.}
{La solución es un número del conjunto de soluciones \(\left \{x\in \mathbb{R} : -1 < x\leq - 5 \right \}\).}{La solución es un número del conjunto de soluciones \(\left \{x\in \mathbb{R} : 5 < x\leq 7 \right \}\).}{La solución es un número del conjunto de soluciones \(\left \{x\in \mathbb{R} : -4 < x\leq - 1 \right \}\).}{La solución es un número del conjunto de soluciones \(\left \{x\in \mathbb{R} : -1 < x\leq 2 \right \}\).}{}{}}
\End

\Data\ID{2010007802}
\Updated{2022-11-01 09:11:50}
\Level{A}
\SubArea{Radical Equations and Inequalities}
\LANGen{
\Question{
Find the domain of the following expression. 
\[ 
 \sqrt{\left (3x - 2 \right ) \left (4+5x\right )}
\]}
{\(\left(-\infty ;-\frac{4}
{5}\right] \cup \left[ \frac{2}
{3};\infty \right )\)}{\(\left[ -\frac{4}
{5}; \frac{2} 
{3}\right] \)}{\(\left (-\infty ;-\frac{4}
{5}\right) \cup \left( \frac{2}
{3};\infty \right)\)}{\(\left( -\frac{4}
{5}; \frac{2} 
{3}\right) \)}{}{}}
\LANGpl{
\Question{
Znajdź dziedzinę wyrażenia. 
\[ 
 \sqrt{\left (3x - 2 \right ) \left (4+5x\right )}
\]}
{\(\left (-\infty ;-\frac{4}
{5}\right\rangle \cup \left\langle \frac{2}
{3};\infty \right )\)}{\(\left\langle -\frac{4}
{5}; \frac{2} 
{3}\right\rangle \)}{\(\left (-\infty ;-\frac{4}
{5}\right) \cup \left( \frac{2}
{3};\infty \right)\)}{\(\left( -\frac{4}
{5}; \frac{2} 
{3}\right) \)}{}{}}
\LANGcs{
\Question{
Jaký je definiční obor tohoto výrazu? 
\[ 
 \sqrt{\left (3x - 2 \right ) \left (4+5x\right )}
\]}
{\(\left (-\infty ;-\frac{4}
{5}\right\rangle \cup \left\langle \frac{2}
{3};\infty \right )\)}{\(\left\langle -\frac{4}
{5}; \frac{2} 
{3}\right\rangle \)}{\(\left (-\infty ;-\frac{4}
{5}\right) \cup \left( \frac{2}
{3};\infty \right)\)}{\(\left( -\frac{4}
{5}; \frac{2} 
{3}\right) \)}{}{}}
\LANGsk{
\Question{
Nájdite definičný obor daného výrazu.
\[ 
 \sqrt{\left (3x - 2 \right ) \left (4+5x\right )}
\]}
{\(\left (-\infty ;-\frac{4}
{5}\right\rangle \cup \left\langle \frac{2}
{3};\infty \right )\)}{\(\left\langle -\frac{4}
{5}; \frac{2} 
{3}\right\rangle \)}{\(\left (-\infty ;-\frac{4}
{5}\right) \cup \left( \frac{2}
{3};\infty \right)\)}{\(\left( -\frac{4}
{5}; \frac{2} 
{3}\right) \)}{}{}}
\LANGes{
\Question{
Halla el dominio de la siguiente expresión.
\[ 
 \sqrt{\left (3x - 2 \right ) \left (4+5x\right )}
\]}
{\(\left (-\infty ;-\frac{4}
{5}\right] \cup \left[ \frac{2}
{3};\infty \right )\)}{\(\left[ -\frac{4}
{5}; \frac{2} 
{3}\right] \)}{\(\left (-\infty ;-\frac{4}
{5}\right) \cup \left( \frac{2}
{3};\infty \right)\)}{\(\left( -\frac{4}
{5}; \frac{2} 
{3}\right) \)}{}{}}
\End

\Data\ID{9000020001}
\Updated{2020-07-15 03:07:34}
\Level{A}
\SubArea{Radical Equations and Inequalities}
\LANGen{
\Question{
Find the domain of the following equation. 
\[ 
 \sqrt{2x - 5} = 3
\]}
{\(\left [ \frac{5}
{2};\infty \right )\)}{\(\left (\frac{2}
{5};\infty \right )\)}{\(\left [ -\frac{5}
{2};\infty \right )\)}{\(\left (\infty ; \frac{2} 
{5}\right )\)}{}{}}
\LANGpl{
\Question{
Znajdź dziedzinę podanego równania. 
\[ 
 \sqrt{2x - 5} = 3
\]}
{\(\left [ \frac{5}
{2};\infty \right )\)}{\(\left (\frac{2}
{5};\infty \right )\)}{\(\left [ -\frac{5}
{2};\infty \right )\)}{\(\left (\infty ; \frac{2} 
{5}\right )\)}{}{}}
\LANGcs{
\Question{
Určete definiční obor rovnice \(\sqrt{2x - 5} = 3\).}
{\(\left \langle \frac{5}
{2};\infty \right )\)}{\(\left (\frac{2}
{5};\infty \right )\)}{\(\left \langle -\frac{5}
{2};\infty \right )\)}{\(\left (\infty ; \frac{2} 
{5}\right )\)}{}{}}
\LANGsk{
\Question{
Určte definičný obor danej rovnice. 
\[ 
 \sqrt{2x - 5} = 3
\]}
{\(\left [ \frac{5}
{2};\infty \right )\)}{\(\left (\frac{2}
{5};\infty \right )\)}{\(\left [ -\frac{5}
{2};\infty \right )\)}{\(\left (\infty ; \frac{2} 
{5}\right )\)}{}{}}
\LANGes{
\Question{
Halla el dominio de la siguiente ecuación. 
\[ 
 \sqrt{2x - 5} = 3
\]}
{\(\left [ \frac{5}
{2};\infty \right )\)}{\(\left (\frac{2}
{5};\infty \right )\)}{\(\left [ -\frac{5}
{2};\infty \right )\)}{\(\left (\infty ; \frac{2} 
{5}\right )\)}{}{}}
\End

\Data\ID{9000020002}
\Updated{2022-04-23 05:04:31}
\Level{A}
\SubArea{Radical Equations and Inequalities}
\LANGen{
\Question{
Find the domain of the following equation. 
\[ 
 \sqrt{6 - x} = 11
\]}
{\((-\infty ;6] \)}{\((5;\infty )\)}{\((-\infty ;5)\)}{\([ - 6;\infty )\)}{}{}}
\LANGpl{
\Question{
Znajdź dziedzinę podanego równania.
\[ 
 \sqrt{6 - x} = 11
\]}
{\((-\infty ;6] \)}{\((5;\infty )\)}{\((-\infty ;5)\)}{\([ - 6;\infty )\)}{}{}}
\LANGcs{
\Question{
Určete definiční obor rovnice \(\sqrt{6 - x} = 11\).}
{\((-\infty ;6\rangle \)}{\((5;\infty )\)}{\((-\infty ;5)\)}{\(\langle - 6;\infty )\)}{}{}}
\LANGsk{
\Question{
Určte definičný obor rovnice. 
\[ 
 \sqrt{6 - x} = 11
\]}
{\((-\infty ;6] \)}{\((5;\infty )\)}{\((-\infty ;5)\)}{\([ - 6;\infty )\)}{}{}}
\LANGes{
\Question{
Halla el dominio de la siguiente ecuación.
\[ 
 \sqrt{6 - x} = 11
\]}
{\((-\infty ;6] \)}{\((5;\infty )\)}{\((-\infty ;5)\)}{\([ - 6;\infty )\)}{}{}}
\End

\Data\ID{9000020003}
\Updated{2022-04-23 05:04:31}
\Level{A}
\SubArea{Radical Equations and Inequalities}
\LANGen{
\Question{
Find the domain of the following equation. 
\[ 
 \sqrt{3x + 6} + \sqrt{8 - 2x} = 11
\]}
{\([ - 2;4] \)}{\((-\infty ;-2] \)}{\([ - 2;\infty )\)}{\([ 4;\infty )\)}{}{}}
\LANGpl{
\Question{
Znajdź dziedzinę podanego równania.
\[ 
 \sqrt{3x + 6} + \sqrt{8 - 2x} = 11
\]}
{\([ - 2;4] \)}{\((-\infty ;-2] \)}{\([ - 2;\infty )\)}{\([ 4;\infty )\)}{}{}}
\LANGcs{
\Question{
Určete definiční obor rovnice \(\sqrt{3x + 6} + \sqrt{8 - 2x} = 11\).}
{\(\langle - 2;4\rangle \)}{\((-\infty ;-2\rangle \)}{\(\langle - 2;\infty )\)}{\(\langle 4;\infty )\)}{}{}}
\LANGsk{
\Question{
Určte definičný obor rovnice. 
\[ 
 \sqrt{3x + 6} + \sqrt{8 - 2x} = 11
\]}
{\([ - 2;4] \)}{\((-\infty ;-2] \)}{\([ - 2;\infty )\)}{\([ 4;\infty )\)}{}{}}
\LANGes{
\Question{
Halla el dominio de la siguiente ecuación. 
\[ 
 \sqrt{3x + 6} + \sqrt{8 - 2x} = 11
\]}
{\([ - 2;4] \)}{\((-\infty ;-2] \)}{\([ - 2;\infty )\)}{\([ 4;\infty )\)}{}{}}
\End

\Data\ID{9000020004}
\Updated{2020-07-15 03:07:34}
\Level{A}
\SubArea{Radical Equations and Inequalities}
\LANGen{
\Question{
Find the domain of the following equation. 
\[ 
 \sqrt{x - 7} + \sqrt{3x + 12} = 5
\]}
{\([ 7;\infty )\)}{\([ - 4;7] \)}{\([ - 4;\infty )\)}{\((-4;7)\)}{}{}}
\LANGpl{
\Question{
Znajdź dziedzinę podanego równania.
\[ 
 \sqrt{x - 7} + \sqrt{3x + 12} = 5
\]}
{\([ 7;\infty )\)}{\([ - 4;7] \)}{\([ - 4;\infty )\)}{\((-4;7)\)}{}{}}
\LANGcs{
\Question{
Určete definiční obor rovnice \(\sqrt{x - 7} + \sqrt{3x + 12} = 5\).}
{\(\langle 7;\infty )\)}{\(\langle - 4;7\rangle \)}{\(\langle - 4;\infty )\)}{\((-4;7)\)}{}{}}
\LANGsk{
\Question{
Určte definičný obor rovnice. 
\[ 
 \sqrt{x - 7} + \sqrt{3x + 12} = 5
\]}
{\([ 7;\infty )\)}{\([ - 4;7] \)}{\([ - 4;\infty )\)}{\((-4;7)\)}{}{}}
\LANGes{
\Question{
Halla el dominio de la siguiente ecuación. 
\[ 
 \sqrt{x - 7} + \sqrt{3x + 12} = 5
\]}
{\([ 7;\infty )\)}{\([ - 4;7] \)}{\([ - 4;\infty )\)}{\((-4;7)\)}{}{}}
\End

\Data\ID{9000020005}
\Updated{2022-04-23 05:04:31}
\Level{A}
\SubArea{Radical Equations and Inequalities}
\LANGen{
\Question{
Identify a true statement referring to the solution of the following equation.
\[ 
 \sqrt{2x - 5} = 3
\]}
{The solution is a prime number.}{The solution is an even number.}{The solution is a number divisible by the number
\(3\).}{The solution is an irrational number.}{}{}}
\LANGpl{
\Question{
Określ, które z poniższych zdań jest prawdziwe w odniesieniu do rozwiązania poniższego równania.
\[ 
 \sqrt{2x - 5} = 3
\]}
{Rozwiązaniem jest liczba pierwsza.}{Rozwiązaniem jest liczba parzysta.}{Rozwiązaniem jest liczba podzielna przez
\(3\).}{Rozwiązaniem jest liczba niewymierna.}{}{}}
\LANGcs{
\Question{
Je dána rovnice \(\sqrt{2x - 5} = 3\).
Vyberte pravdivé tvrzení.}
{Kořenem rovnice je prvočíslo.}{Kořenem rovnice je sudé číslo.}{Kořen rovnice je dělitelný třemi.}{Kořenem rovnice je iracionální číslo.}{}{}}
\LANGsk{
\Question{
Je daná rovnica. Vyberte pravdivé tvrdenie.
\[ 
 \sqrt{2x - 5} = 3
\]}
{Koreňom rovnice je prvočíslo.}{Koreňom rovnice je párne číslo.}{Koreň rovnice je deliteľný číslom
\(3\).}{Koreňom rovnice je iracionálne číslo.}{}{}}
\LANGes{
\Question{
Identifica la proposición lógica que se refiera a la solución de la siguiente ecuación:
\[ 
 \sqrt{2x - 5} = 3
\]}
{La solución es un número primo.}{La solución es un número par.}{La solución es un número divisible por \(3\).}{La solución es un número irracional.}{}{}}
\End

\Data\ID{9000020006}
\Updated{2020-12-29 08:12:11}
\Level{A}
\SubArea{Radical Equations and Inequalities}
\LANGen{
\Question{
Identify a true statement referring to the solution of the following equation.
\[ 
 \sqrt{3x - 8} = x - 6
\]}
{The equation has a unique solution and this solution is an odd number.}{The equation has two solutions, the sum of these solutions is divisible by
\(5\).}{The equation has a unique solution and this solution is an even number.}{The equation does not have a solution in
\(\mathbb{R}\).}{}{}}
\LANGpl{
\Question{
Wskaż, które z poniższych zdań jest prawdziwe w odniesieniu do rozwiązania podanego równania.
\[ 
 \sqrt{3x - 8} = x - 6
\]}
{Równanie ma tylko jedno rozwiązanie i jest nim liczba nieparzysta.}{Równanie ma dwa rozwiązania, których suma jest podzielna przez
\(5\).}{Równanie ma tylko jedno rozwiązanie i jest nim liczba parzysta.}{Równanie nie ma rozwiązania w przedziale
\(\mathbb{R}\).}{}{}}
\LANGcs{
\Question{
Je dána rovnice \(\sqrt{3x - 8} = x - 6\).
Vyberte pravdivé tvrzení.}
{Rovnice má právě jeden kořen a je to liché číslo.}{Rovnice má dva kořeny, jejichž součet je násobkem pěti.}{Rovnice má právě jeden kořen a je to sudé číslo.}{Rovnice nemá v \(\mathbb{R}\) 
řešení.}{}{}}
\LANGsk{
\Question{
Je daná rovnica. Vyberte pravdivé tvrdenie.
\[ 
 \sqrt{3x - 8} = x - 6
\]}
{Rovnica má práve jeden koreň a koreňom rovnice je nepárne číslo.}{Rovnica má dva korene, ktorých súčet je deliteľný 
\(5\).}{Rovnica má práve jeden koreň a koreňom rovnice je párne číslo.}{Rovnica nemá riešenie v 
\(\mathbb{R}\).}{}{}}
\LANGes{
\Question{
Dada la ecuación:
\[ 
 \sqrt{3x - 8} = x - 6
\]
Identifica la proposición lógica que hace referencia a la ecuación.}
{La ecuación tiene solo una solución, es un número impar.}{La ecuación tiene dos soluciones, la suma de estas soluciones es divisible por \(5\).}{La ecuación tiene solo una solución, es un número par.}{La ecuación no tiene soluciones en \(\mathbb{R}\).}{}{}}
\End

\Data\ID{9000020007}
\Updated{2020-12-29 08:12:35}
\Level{A}
\SubArea{Radical Equations and Inequalities}
\LANGen{
\Question{
Identify a true statement referring to the solution of the following equation.
\[ 
 \sqrt{x^{2 } - 4} = x + 1
\]}
{The equation does not have a solution in
\(\mathbb{R}\).}{The equation has a unique negative solution.}{The equation has a unique positive solution.}{The equation has two solutions.}{}{}}
\LANGpl{
\Question{
Które z poniższych zdań jest prawdziwe w odniesieniu do rozwiązania podanego równania?
\[ 
 \sqrt{x^{2 } - 4} = x + 1
\]}
{Równanie nie ma rozwiązania w przedziale
\(\mathbb{R}\).}{Równanie ma jedno rozwiązanie i jest nim liczba ujemna.}{Równanie ma jedno rozwiązanie i jest nim liczba dodatnia.}{Równanie ma dwa rozwiązania.}{}{}}
\LANGcs{
\Question{
Je dána rovnice \(\sqrt{x^{2 } - 4} = x + 1\).
Vyberte pravdivé tvrzení.}
{Rovnice nemá v \(\mathbb{R}\) 
řešení.}{Rovnice má právě jeden záporný kořen.}{Rovnice má právě jeden kladný kořen.}{Rovnice má dva kořeny.}{}{}}
\LANGsk{
\Question{
Je daná rovnica. Vyberte pravdivé tvrdenie.
\[ 
 \sqrt{x^{2 } - 4} = x + 1
\]}
{Rovnica nemá riešenie v
\(\mathbb{R}\).}{Rovnica má práve jeden záporný koreň.}{Rovnica má práve jeden kladný koreň.}{Rovnica má dva korene.}{}{}}
\LANGes{
\Question{
Dada la ecuación:
\[ 
 \sqrt{x^{2 } - 4} = x + 1
\]
Identifica la proposición lógica que hace referencia a la ecuación.}
{La ecuación no tiene soluciones en \(\mathbb{R}\).}{La ecuación tiene unicamente una solución negativa.}{La ecuación tiene unicamente una solución positiva.}{La ecuación tiene dos soluciones.}{}{}}
\End

\Data\ID{9000020008}
\Updated{2020-12-29 08:12:13}
\Level{A}
\SubArea{Radical Equations and Inequalities}
\LANGen{
\Question{
Identify a true statement referring to the solution of the following equation.
\[ 
 6x - 13\sqrt{x} + 6 = 0
\]
 Hint: Use the substitution \(y = \sqrt{x}\).}
{The solutions \(x_{1}\) 
and \(x_{2}\) 
satisfy \(x_{1} = \frac{1} 
{x_{2}} \).}{The equation has a unique solution
\(x_{1}\). This solution
satisfies \(x_{1} < 1\).}{The equation has a unique solution
\(x_{1}\). This solution
satisfies \(x_{1} > 1\).}{The equation does not have a solution in
\(\mathbb{R}\).}{}{}}
\LANGpl{
\Question{
Określ, które z poniższych zdań jest prawdziwe w odniesieniu do rozwiązania podanego równania.
\[ 
 6x - 13\sqrt{x} + 6 = 0
\]
 Wskazówka: Użyj  substytucja \(y = \sqrt{x}\).}
{Rozwiązania \(x_{1}\) 
i \(x_{2}\) 
spełniają \(x_{1} = \frac{1} 
{x_{2}} \).}{Równanie ma tylko jedno rozwiązanie
\(x_{1}\). To rozwiązanie spełnia \(x_{1} < 1\).}{Równanie ma tylko jedno rozwiązanie
\(x_{1}\). To rozwiązanie spełnia 
 \(x_{1} > 1\).}{Równanie nie ma rozwiązania w przedziale
\(\mathbb{R}\).}{}{}}
\LANGcs{
\Question{
Je dána rovnice \(6x - 13\sqrt{x} + 6 = 0\).
Vyberte pravdivé tvrzení. Nápověda: Využijte substituce \(y = \sqrt{x}\).}
{Rovnice má kořeny \(x_{1}\) 
a \(x_{2}\),
\(x_{1} = \frac{1} 
{x_{2}} \).}{Rovnice má právě jeden kořen
\(x_{1}\) takový, že
\(x_{1} < 1\).}{Rovnice má právě jeden kořen
\(x_{1}\) takový, že
\(x_{1} > 1\).}{Rovnice nemá v \(\mathbb{R}\) 
řešení.}{}{}}
\LANGsk{
\Question{
Je daná rovnica. Vyberte pravdivé tvrdenie.
\[ 
 6x - 13\sqrt{x} + 6 = 0
\]
 Využite substitúciu \(y = \sqrt{x}\).}
{Rovnica má korene \(x_{1}\) 
a \(x_{2}\),
 \(x_{1} = \frac{1} 
{x_{2}} \).}{Rovnica má práve jeden koreň
\(x_{1}\), \(x_{1} < 1\).}{Rovnica má práve jeden koreň
\(x_{1}\), 
 \(x_{1} > 1\).}{Rovnica nemá riešenie v
\(\mathbb{R}\).}{}{}}
\LANGes{
\Question{
Dada la ecuación: 
\[ 
 6x - 13\sqrt{x} + 6 = 0
\]
Identifica la proposición lógica que hace referencia a la ecuación.
Sugerencia: usa la sustitución \(y = \sqrt{x}\).}
{Las soluciones \(x_{1}\) 
y \(x_{2}\) 
satisfacen \(x_{1} = \frac{1} 
{x_{2}} \).}{La ecuación tiene solo una solución
\(x_{1}\). La solución satisface \(x_{1} < 1\).}{La ecuación tiene solo una solución
\(x_{1}\). La solución satisface \(x_{1} > 1\).}{La ecuación no tiene soluciones en \(\mathbb{R}\).}{}{}}
\End

\Data\ID{9000020009}
\Updated{2020-07-15 03:07:34}
\Level{A}
\SubArea{Radical Equations and Inequalities}
\LANGen{
\Question{
Choose the equation which is obtained by squaring both
sides of the following equation. 
\[ 
 \sqrt{3x + 2} = x - 6
\]}
{\(x^{2} - 15x + 34 = 0\)}{\(x^{2} - 3x - 38 = 0\)}{\(x^{2} - 3x - 34 = 0\)}{\(x^{2} - 15x - 38 = 0\)}{}{}}
\LANGpl{
\Question{
Wybierz równanie otrzymanego przez podniesienie do kwadratu obu stron następującego równania.
\[ 
 \sqrt{3x + 2} = x - 6
\]}
{\(x^{2} - 15x + 34 = 0\)}{\(x^{2} - 3x - 38 = 0\)}{\(x^{2} - 3x - 34 = 0\)}{\(x^{2} - 15x - 38 = 0\)}{}{}}
\LANGcs{
\Question{
Je dána rovnice \(\sqrt{3x + 2} = x - 6\).
Vyberte rovnici, kterou získáte po umocnění dané rovnice na druhou.}
{\(x^{2} - 15x + 34 = 0\)}{\(x^{2} - 3x - 38 = 0\)}{\(x^{2} - 3x - 34 = 0\)}{\(x^{2} - 15x - 38 = 0\)}{}{}}
\LANGsk{
\Question{
Vyberte rovnicu, ktorú získate po umocnení danej rovnice na druhú. 
\[ 
 \sqrt{3x + 2} = x - 6
\]}
{\(x^{2} - 15x + 34 = 0\)}{\(x^{2} - 3x - 38 = 0\)}{\(x^{2} - 3x - 34 = 0\)}{\(x^{2} - 15x - 38 = 0\)}{}{}}
\LANGes{
\Question{
Elige la ecuación que se obtiene al elevar al cuadrado ambos lados de la siguiente ecuación.
\[ 
 \sqrt{3x + 2} = x - 6
\]}
{\(x^{2} - 15x + 34 = 0\)}{\(x^{2} - 3x - 38 = 0\)}{\(x^{2} - 3x - 34 = 0\)}{\(x^{2} - 15x - 38 = 0\)}{}{}}
\End

\Data\ID{9000020010}
\Updated{2020-07-15 03:07:34}
\Level{A}
\SubArea{Radical Equations and Inequalities}
\LANGen{
\Question{
Choose the equation which is obtained by squaring both
sides of the following equation. 
\[ 
 \sqrt{x^{2 } - x + 5} = 2x - 5
\]}
{\(3x^{2} - 19x + 20 = 0\)}{\(x^{2} + 3x + 20 = 0\)}{\(3x^{2} + x - 30 = 0\)}{\(3x^{2} + x + 20 = 0\)}{}{}}
\LANGpl{
\Question{
Wybierz równanie otrzymanego przez podniesienie do kwadratu obu stron następującego równania.
\[ 
 \sqrt{x^{2 } - x + 5} = 2x - 5
\]}
{\(3x^{2} - 19x + 20 = 0\)}{\(x^{2} + 3x + 20 = 0\)}{\(3x^{2} + x - 30 = 0\)}{\(3x^{2} + x + 20 = 0\)}{}{}}
\LANGcs{
\Question{
Je dána rovnice \(\sqrt{x^{2 } - x + 5} = 2x - 5\).
Vyberte rovnici, kterou získáte po umocnění dané rovnice na druhou.}
{\(3x^{2} - 19x + 20 = 0\)}{\(x^{2} + 3x + 20 = 0\)}{\(3x^{2} + x - 30 = 0\)}{\(3x^{2} + x + 20 = 0\)}{}{}}
\LANGsk{
\Question{
Vyberte rovnicu, ktorú získate po umocnení danej rovnice na druhú. 
\[ 
 \sqrt{x^{2 } - x + 5} = 2x - 5
\]}
{\(3x^{2} - 19x + 20 = 0\)}{\(x^{2} + 3x + 20 = 0\)}{\(3x^{2} + x - 30 = 0\)}{\(3x^{2} + x + 20 = 0\)}{}{}}
\LANGes{
\Question{
Elige la ecuación que se obtiene al elevar al cuadrado ambos lados de la siguiente ecuación.
\[ 
 \sqrt{x^{2 } - x + 5} = 2x - 5
\]}
{\(3x^{2} - 19x + 20 = 0\)}{\(x^{2} + 3x + 20 = 0\)}{\(3x^{2} + x - 30 = 0\)}{\(3x^{2} + x + 20 = 0\)}{}{}}
\End

\Data\ID{9000022305}
\Updated{2020-07-15 03:07:34}
\Level{A}
\SubArea{Radical Equations and Inequalities}
\LANGen{
\Question{
Find the domain of the following expression. 
\[ 
 \sqrt{-x^{2 } + 16x - 63}
\]}
{\(\left [ 7;9\right ] \)}{\(\left (-\infty ;7\right )\cup \left (9;\infty \right )\)}{\(\left (-\infty ;-7\right ] \cup \left [ 9;\infty \right )\)}{\(\left (7;9\right )\)}{\(\left [ -7;9\right ] \)}{}}
\LANGpl{
\Question{
Znajdź dziedzinę podanego wyrażenia.
\[ 
 \sqrt{-x^{2 } + 16x - 63}
\]}
{\(\left [ 7;9\right ] \)}{\(\left (-\infty ;7\right )\cup \left (9;\infty \right )\)}{\(\left (-\infty ;-7\right ] \cup \left [ 9;\infty \right )\)}{\(\left (7;9\right )\)}{\(\left [ -7;9\right ] \)}{}}
\LANGcs{
\Question{
Výraz \(\sqrt{-x^{2 } + 16x - 63}\) 
má definiční obor:}
{\(\left \langle 7;9\right \rangle \)}{\(\left (-\infty ;7\right )\cup \left (9;\infty \right )\)}{\(\left (-\infty ;-7\right \rangle \cup \left \langle 9;\infty \right )\)}{\(\left (7;9\right )\)}{\(\left \langle -7;9\right \rangle \)}{}}
\LANGsk{
\Question{
Určte definičný obor daného výrazu.
\[ 
 \sqrt{-x^{2 } + 16x - 63}
\]}
{\(\left [ 7;9\right ] \)}{\(\left (-\infty ;7\right )\cup \left (9;\infty \right )\)}{\(\left (-\infty ;-7\right ] \cup \left [ 9;\infty \right )\)}{\(\left (7;9\right )\)}{\(\left [ -7;9\right ] \)}{}}
\LANGes{
\Question{
Halla el dominio de la siguiente expresión. 
\[ 
 \sqrt{-x^{2 } + 16x - 63}
\]}
{\(\left [ 7;9\right ] \)}{\(\left (-\infty ;7\right )\cup \left (9;\infty \right )\)}{\(\left (-\infty ;-7\right ] \cup \left [ 9;\infty \right )\)}{\(\left (7;9\right )\)}{\(\left [ -7;9\right ] \)}{}}
\End

\Data\ID{9000022801}
\Updated{2020-07-15 03:07:34}
\Level{A}
\SubArea{Radical Equations and Inequalities}
\LANGen{
\Question{
Find the domain of the expression. 
\[ 
 \sqrt{x^{2 } + x - 2}
\]}
{\(\left (-\infty ;-2\right ] \cup \left [ 1;\infty \right )\)}{\(\left [ -2;1\right ] \)}{\(\left (-2;1\right )\)}{\(\left (-\infty ;-2\right )\cup \left (1;\infty \right )\)}{}{}}
\LANGpl{
\Question{
Znajdź dziedzinę podanego wyrażenia.
\[ 
 \sqrt{x^{2 } + x - 2}
\]}
{\(\left (-\infty ;-2\right ] \cup \left [ 1;\infty \right )\)}{\(\left [ -2;1\right ] \)}{\(\left (-2;1\right )\)}{\(\left (-\infty ;-2\right )\cup \left (1;\infty \right )\)}{}{}}
\LANGcs{
\Question{
Množina všech \(x\in \mathbb{R}\),
pro která je výraz \(\sqrt{x^{2 } + x - 2}\) 
definován, je:}
{\(\left (-\infty ;-2\right \rangle \cup \left \langle 1;\infty \right )\)}{\(\left \langle -2;1\right \rangle \)}{\(\left (-2;1\right )\)}{\(\left (-\infty ;-2\right )\cup \left (1;\infty \right )\)}{}{}}
\LANGsk{
\Question{
Určte definičný obor daného výrazu. 
\[ 
 \sqrt{x^{2 } + x - 2}
\]}
{\(\left (-\infty ;-2\right ] \cup \left [ 1;\infty \right )\)}{\(\left [ -2;1\right ] \)}{\(\left (-2;1\right )\)}{\(\left (-\infty ;-2\right )\cup \left (1;\infty \right )\)}{}{}}
\LANGes{
\Question{
Halla el dominio de la siguiente expresión.
\[ 
 \sqrt{x^{2 } + x - 2}
\]}
{\(\left (-\infty ;-2\right ] \cup \left [ 1;\infty \right )\)}{\(\left [ -2;1\right ] \)}{\(\left (-2;1\right )\)}{\(\left (-\infty ;-2\right )\cup \left (1;\infty \right )\)}{}{}}
\End

\Data\ID{9000023701}
\Updated{2020-12-29 08:12:18}
\Level{A}
\SubArea{Radical Equations and Inequalities}
\LANGen{
\Question{
Identify a true statement which concerns the following equation. 
\[ 
 \sqrt{x - 3} = 1
\]}
{The solution is \(x = 4\).}{The solution is \(x = 2\).}{The solution is \(x = 5\).}{The equation does not have any solution.}{}{}}
\LANGpl{
\Question{
Które z poniższych zdań jest prawdziwe w odniesieniu do podanego równania?
\[ 
 \sqrt{x - 3} = 1
\]}
{Rozwiązaniem jest \(x = 4\).}{Rozwiązaniem jest \(x = 2\).}{Rozwiązaniem jest \(x = 5\).}{Równanie nie ma rozwiązania.}{}{}}
\LANGcs{
\Question{
Je dána rovnice \(\sqrt{x - 3} = 1\).
Které z následujících tvrzení je správné?}
{Řešením této rovnice je číslo
\(4\).}{Řešením této rovnice je číslo
\(2\).}{Řešením této rovnice je číslo
\(5\).}{Tato rovnice nemá řešení.}{}{}}
\LANGsk{
\Question{
Je daná rovnica. Ktoré z nasledujúcich tvrdení je správne? 
\[ 
 \sqrt{x - 3} = 1
\]}
{Koreňom tejto rovnice je \(x = 4\).}{Koreňom tejto rovnice je \(x = 2\).}{Koreňom tejto rovnice je \(x = 5\).}{Rovnica nemá žiadne riešenie.}{}{}}
\LANGes{
\Question{
Dada la ecuación:
\[ 
 \sqrt{x - 3} = 1
\]
Identifica la proposición lógica que hace referencia a la ecuación.}
{La solución es \(x = 4\).}{La solución es \(x = 2\).}{La solución es \(x = 5\).}{La ecuación no tiene solución.}{}{}}
\End

\Data\ID{9000023702}
\Updated{2020-12-29 08:12:11}
\Level{A}
\SubArea{Radical Equations and Inequalities}
\LANGen{
\Question{
Identify a true statement which concerns the following equation. 
\[ 
 \sqrt{x + 7} = 3
\]}
{The solution is \(x = 2\).}{The solution is \(x = -4\).}{The solution is \(x = -2\).}{The equation does not have a solution.}{}{}}
\LANGpl{
\Question{
Które z poniższych zdań jest prawdziwe w odniesieniu do następującego równania?
\[ 
 \sqrt{x + 7} = 3
\]}
{Rozwiązaniem jest  \(x = 2\).}{Rozwiązaniem jest \(x = -4\).}{Rozwiązaniem jest \(x = -2\).}{Równanie nie ma rozwiązania.}{}{}}
\LANGcs{
\Question{
Je dána rovnice \(\sqrt{x + 7} = 3\).
Které z následujících tvrzení je správné?}
{Řešením této rovnice je číslo
\(2\).}{Řešením této rovnice je číslo
\(- 4\).}{Řešením této rovnice je číslo
\(- 2\).}{Tato rovnice nemá řešení.}{}{}}
\LANGsk{
\Question{
Je daná rovnica. Ktoré z nasledujúcich tvrdení je správne? 
\[ 
 \sqrt{x + 7} = 3
\]}
{Koreň rovnice je \(x = 2\).}{Koreň rovnice je \(x = -4\).}{Koreň rovnice je \(x = -2\).}{Táto rovnica nemá riešenie.}{}{}}
\LANGes{
\Question{
Dada la ecuación:
\[ 
 \sqrt{x + 7} = 3
\]
Identifica la proposición lógica que hace referencia a la ecuación.}
{La solución es \(x = 2\).}{La solución es \(x = -4\).}{La solución es \(x = -2\).}{La ecuación no tiene solución.}{}{}}
\End

\Data\ID{9000023703}
\Updated{2020-12-29 08:12:40}
\Level{A}
\SubArea{Radical Equations and Inequalities}
\LANGen{
\Question{
Identify a true statement which concerns the following equation. 
\[ 
 \sqrt{x + 1} = 2
\]}
{The solution is a number from the interval
\([ 2;5)\).}{The solution is a number from the interval
\([ - 1;2] \).}{The solution is a number from the interval
\([ - 2;3)\).}{The solution is a number from the interval
\((4;7)\).}{}{}}
\LANGpl{
\Question{
Które z poniższych zdań jest prawdziwe w odniesieniu do rozwiązania podanego równania?
\[ 
 \sqrt{x + 1} = 2
\]}
{Rozwiązaniem jest liczba mieszcząca się w przedziale
\([ 2;5)\).}{Rozwiązaniem jest liczba mieszcząca się w przedziale
\([ - 1;2] \).}{Rozwiązaniem jest liczba mieszcząca się w przedziale
\([ - 2;3)\).}{Rozwiązaniem jest liczba mieszcząca się w przedziale
\((4;7)\).}{}{}}
\LANGcs{
\Question{
Je dána rovnice \(\sqrt{x + 1} = 2\).
Které z následujících tvrzení je správné?}
{Řešením této rovnice je číslo z intervalu
\(\langle 2;5)\).}{Řešením této rovnice je číslo z intervalu
\(\langle - 1;2\rangle \).}{Řešením této rovnice je číslo z intervalu
\(\langle - 2;3)\).}{Řešením této rovnice je číslo z intervalu
\((4;7)\).}{}{}}
\LANGsk{
\Question{
Je daná rovnica. Ktoré z nasledujúcich tvrdení je správne? 
\[ 
 \sqrt{x + 1} = 2
\]}
{Riešením tejto rovnice je číslo z intervalu
\([ 2;5)\).}{Riešením tejto rovnice je číslo z intervalu
\([ - 1;2] \).}{Riešením tejto rovnice je číslo z  intervalu
\([ - 2;3)\).}{Riešením tejto rovnice je číslo z intervalu
\((4;7)\).}{}{}}
\LANGes{
\Question{
Dada la ecuación:
\[ 
 \sqrt{x + 1} = 2
\]
Identifica la proposición lógica que hace referencia a la ecuación.}
{La solución es un número del intervalo \([ 2;5)\).}{La solución es un número del intervalo \([ - 1;2] \).}{La solución es un número del intervalo \([ - 2;3)\).}{La solución es un número del intervalo \((4;7)\).}{}{}}
\End

\Data\ID{9000023704}
\Updated{2020-12-29 08:12:52}
\Level{A}
\SubArea{Radical Equations and Inequalities}
\LANGen{
\Question{
Identify a true statement which concerns the following equation. 
\[ 
 \sqrt{x + 20} = 4
\]}
{The solution is from the set \(B = \left \{x\in \mathbb{R} : -6\leq x\leq - 2\right \}\).}{The solution is from the set \(A = \left \{x\in \mathbb{R} : -4 < x\leq - 1\right \}\).}{The solution is from the set \(C = \left \{x\in \mathbb{R} : -7\leq x\leq - 5\right \}\).}{The solution is from the set \(D = \left \{x\in \mathbb{R} : -3 < x < 0\right \}\).}{}{}}
\LANGpl{
\Question{
Które z poniższych zdań jest prawdziwe w odniesieniu do rozwiązania podanego równania?
\[ 
 \sqrt{x + 20} = 4
\]}
{Rozwiązanie należy do zbioru \(B = \left \{x\in \mathbb{R} : -6\leq x\leq - 2\right \}\).}{Rozwiązanie należy do zbioru  \(A = \left \{x\in \mathbb{R} : -4 < x\leq - 1\right \}\).}{Rozwiązanie należy do zbioru \(C = \left \{x\in \mathbb{R} : -7\leq x\leq - 5\right \}\).}{Rozwiązanie należy do zbioru  \(D = \left \{x\in \mathbb{R} : -3 < x < 0\right \}\).}{}{}}
\LANGcs{
\Question{
Je dána rovnice \(\sqrt{x + 20} = 4\).
Které z následujících tvrzení je správné?}
{Řešením této rovnice je číslo z množiny
\(B = \left \{x\in \mathbb{R} : -6\leq x\leq - 2\right \}\).}{Řešením této rovnice je číslo z množiny
\(A = \left \{x\in \mathbb{R} : -4 < x\leq - 1\right \}\).}{Řešením této rovnice je číslo z množiny
\(C = \left \{x\in \mathbb{R} : -7\leq x\leq - 5\right \}\).}{Řešením této rovnice je číslo z množiny
\(D = \left \{x\in \mathbb{R} : -3 < x < 0\right \}\).}{}{}}
\LANGsk{
\Question{
Je daná rovnica. Ktoré z nasledujúcich tvrdení je správne? 
\[ 
 \sqrt{x + 20} = 4
\]}
{Riešením tejto rovnice je číslo z množiny \(B = \left \{x\in \mathbb{R} : -6\leq x\leq - 2\right \}\).}{Riešením tejto rovnice je číslo z množiny \(A = \left \{x\in \mathbb{R} : -4 < x\leq - 1\right \}\).}{Riešením tejto rovnice je číslo z množiny \(C = \left \{x\in \mathbb{R} : -7\leq x\leq - 5\right \}\).}{Riešením tejto rovnice je číslo z množiny \(D = \left \{x\in \mathbb{R} : -3 < x < 0\right \}\).}{}{}}
\LANGes{
\Question{
Dada la ecuación:
\[ 
 \sqrt{x + 20} = 4
\]
Identifica la proposición lógica que hace referencia a la ecuación.}
{La solución es del conjunto de soluciones \(B = \left \{x\in \mathbb{R} : -6\leq x\leq - 2\right \}\).}{La solución es del conjunto de soluciones \(A = \left \{x\in \mathbb{R} : -4 < x\leq - 1\right \}\).}{La solución es del conjunto de soluciones \(C = \left \{x\in \mathbb{R} : -7\leq x\leq - 5\right \}\).}{La solución es del conjunto de soluciones \(D = \left \{x\in \mathbb{R} : -3 < x < 0\right \}\).}{}{}}
\End

\Data\ID{9000023705}
\Updated{2020-12-29 08:12:24}
\Level{A}
\SubArea{Radical Equations and Inequalities}
\LANGen{
\Question{
Identify a true statement which concerns the following equation. 
\[ 
 \sqrt{x + 4} = 3
\]}
{The solution is a divisor of \(20\).}{The solution is a divisor of \(6\).}{The solution is a divisor of \(12\).}{The solution is a divisor of \(18\).}{}{}}
\LANGpl{
\Question{
Które z poniższych zdań jest prawdziwe w odniesieniu do podanego równania?
\[ 
 \sqrt{x + 4} = 3
\]}
{Rozwiązaniem jest dzielnik liczby \(20\).}{Rozwiązaniem jest dzielnik liczby  \(6\).}{Rozwiązaniem jest dzielnik liczby  \(12\).}{Rozwiązaniem jest dzielnik liczby \(18\).}{}{}}
\LANGcs{
\Question{
Je dána rovnice \(\sqrt{x + 4} = 3\).
Které z následujících tvrzení je správné?}
{Řešením této rovnice je číslo, které je dělitelem čísla
\(20\).}{Řešením této rovnice je číslo, které je dělitelem čísla
\(6\).}{Řešením této rovnice je číslo, které je dělitelem čísla
\(12\).}{Řešením této rovnice je číslo, které je dělitelem čísla
\(18\).}{}{}}
\LANGsk{
\Question{
Je daná rovnica. Ktoré z nasledujúcich tvrdení je správne? 
\[ 
 \sqrt{x + 4} = 3
\]}
{Riešením tejto rovnice je číslo, ktoré je deliteľom čísla \(20\).}{Riešením tejto rovnice je číslo, ktoré je deliteľom čísla \(6\).}{Riešením tejto rovnice je číslo, ktoré je deliteľom čísla \(12\).}{Riešením tejto rovnice je číslo, ktoré je deliteľom čísla \(18\).}{}{}}
\LANGes{
\Question{
Dada la ecuación:
\[ 
 \sqrt{x + 4} = 3
\]
Identifica la proposición lógica que hace referencia a la ecuación.}
{La solución es un divisor de \(20\).}{La solución es un divisor de \(6\).}{La solución es un divisor de \(12\).}{La solución es un divisor de \(18\).}{}{}}
\End

\Data\ID{9000023706}
\Updated{2020-12-29 08:12:44}
\Level{A}
\SubArea{Radical Equations and Inequalities}
\LANGen{
\Question{
Identify a true statement which concerns the following equation. 
\[ 
 \sqrt{2x + 7} = 5
\]}
{The solution is a multiple of \(3\).}{The solution is a multiple of \(2\).}{The solution is a multiple of \(4\).}{The solution is a multiple of \(5\).}{}{}}
\LANGpl{
\Question{
Które z poniższych stwierdzeń dotyczących danej funkcji jest prawdziwe?
\[ 
 \sqrt{2x + 7} = 5
\]}
{Rozwiązaniem jest wielokrotność liczby  \(3\).}{Rozwiązaniem jest wielokrotność liczby \(2\).}{Rozwiązaniem jest wielokrotność liczby \(4\).}{Rozwiązaniem jest wielokrotność liczby  \(5\).}{}{}}
\LANGcs{
\Question{
Je dána rovnice \(\sqrt{2x + 7} = 5\).
Které z následujících tvrzení je správné?}
{Řešením této rovnice je číslo, které je násobkem čísla
\(3\).}{Řešením této rovnice je číslo, které je násobkem čísla
\(2\).}{Řešením této rovnice je číslo, které je násobkem čísla
\(4\).}{Řešením této rovnice je číslo, které je násobkem čísla
\(5\).}{}{}}
\LANGsk{
\Question{
Je daná rovnica. Ktoré z nasledujúcich tvrdení je správne? 
\[ 
 \sqrt{2x + 7} = 5
\]}
{Riešením tejto rovnice je číslo, ktoré je násobkom čísla \(3\).}{Riešením tejto rovnice je číslo, ktoré je násobkom čísla \(2\).}{Riešením tejto rovnice je číslo, ktoré je násobkom čísla \(4\).}{Riešením tejto rovnice je číslo, ktoré je násobkom čísla \(5\).}{}{}}
\LANGes{
\Question{
Dada la ecuación:
\[ 
 \sqrt{2x + 7} = 5
\]
Identifica la proposición lógica que hace referencia a la ecuación.}
{La solución es un múltiplo de \(3\).}{La solución es un múltiplo de \(2\).}{La solución es un múltiplo de \(4\).}{La solución es un múltiplo de \(5\).}{}{}}
\End

\Data\ID{9000023707}
\Updated{2020-12-29 08:12:48}
\Level{A}
\SubArea{Radical Equations and Inequalities}
\LANGen{
\Question{
Identify a true statement which concerns the following equation. 
\[ 
 \sqrt{3x - 5} = 4
\]}
{The solution is a prime number.}{The solution is from the interval \([ - 5;5] \).}{The solution is from the set \(A = \left \{x\in \mathbb{R} : -4 < x\leq 3\right \}\).}{The solution is a multiple of \(4\).}{}{}}
\LANGpl{
\Question{
Które z poniższych stwierdzeń dotyczących danej funkcji jest prawdziwe?
\[ 
 \sqrt{3x - 5} = 4
\]}
{Rozwiązaniem jest liczba pierwsza.}{Rozwiązanie mieści się w przedziale  \([ - 5;5] \).}{Rozwiązanie należy do zbioru  \(A = \left \{x\in \mathbb{R} : -4 < x\leq 3\right \}\).}{Rozwiązaniem jest wielokrotność liczby \(4\).}{}{}}
\LANGcs{
\Question{
Je dána rovnice \(\sqrt{3x - 5} = 4\).
Které z následujících tvrzení je správné?}
{Řešením této rovnice je číslo, které je prvočíslo.}{Řešením této rovnice je číslo z intervalu
\(\langle - 5;5\rangle \).}{Řešením této rovnice je číslo z množiny
\(A = \left \{x\in \mathbb{R} : -4 < x\leq 3\right \}\).}{Řešením této rovnice je číslo, které je násobkem čísla
\(4\).}{}{}}
\LANGsk{
\Question{
Je daná rovnica. Ktoré z nasledujúcich tvrdení je správne? 
\[ 
 \sqrt{3x - 5} = 4
\]}
{Riešením tejto rovnice je prvočíslo.}{Riešením tejto rovnice je číslo z intervalu \([ - 5;5] \).}{Riešením tejto rovnice je číslo z množiny \(A = \left \{x\in \mathbb{R} : -4 < x\leq 3\right \}\).}{Riešením tejto rovnice je číslo, ktoré je násobkom čísla \(4\).}{}{}}
\LANGes{
\Question{
Dada la ecuación:
\[ 
 \sqrt{3x - 5} = 4
\]
Identifica la proposición lógica que hace referencia a la ecuación.}
{La solución es un número primo.}{La solución es un número del intervalo \([ - 5;5] \).}{La solución es un número del conjunto de soluciones \(A = \left \{x\in \mathbb{R} : -4 < x\leq 3\right \}\).}{La solución es un múltiplo de \(4\).}{}{}}
\End

\Data\ID{9000023708}
\Updated{2020-12-29 09:12:10}
\Level{A}
\SubArea{Radical Equations and Inequalities}
\LANGen{
\Question{
Identify a true statement which concerns the following equation. 
\[ 
 \sqrt{x + 5} = x - 1
\]}
{The solution is an even number.}{The solution is from the interval \([ - 2;2)\).}{The solution is from the set \(A = \left \{x\in \mathbb{R} : -1\leq x < 3\right \}\).}{The solution is a divisor of \(6\).}{}{}}
\LANGpl{
\Question{
Które z poniższych stwierdzeń dotyczących danej funkcji jest prawdziwe?
\[ 
 \sqrt{x + 5} = x - 1
\]}
{Rozwiązaniem jest liczba parzysta.}{Rozwiązanie mieści się w przedziale  \([ - 2;2)\).}{Rozwiązanie należy do zbioru \(A = \left \{x\in \mathbb{R} : -1\leq x < 3\right \}\).}{Rozwiązaniem jest dzielnik liczby \(6\).}{}{}}
\LANGcs{
\Question{
Je dána rovnice \(\sqrt{x + 5} = x - 1\).
Které z následujících tvrzení je správné?}
{Řešením této rovnice je číslo, které je sudé.}{Řešením této rovnice je číslo z intervalu
\(\langle - 2;2)\).}{Řešením této rovnice je číslo z množiny
\(A = \left \{x\in \mathbb{R} : -1\leq x < 3\right \}\).}{Řešením této rovnice je číslo, které je dělitelem čísla
\(6\).}{}{}}
\LANGsk{
\Question{
Je daná rovnica. Ktoré z nasledujúcich tvrdení je správne? 
\[ 
 \sqrt{x + 5} = x - 1
\]}
{Riešením tejto rovnice je číslo, ktoré je párne.}{Riešením tejto rovnice je číslo z intervalu \([ - 2;2)\).}{Riešením tejto rovnice je číslo z množiny \(A = \left \{x\in \mathbb{R} : -1\leq x < 3\right \}\).}{Riešením tejto rovnice je číslo, ktoré je deliteľom čísla  \(6\).}{}{}}
\LANGes{
\Question{
Dada la ecuación:
\[ 
 \sqrt{x + 5} = x - 1
\]
Identifica la proposición lógica que hace referencia a la ecuación.}
{La solución es un número par.}{La solución es un número del intervalo \([ - 2;2)\).}{La solución es un número del conjunto de soluciones  \(A = \left \{x\in \mathbb{R} : -1\leq x < 3\right \}\).}{La solución es un divisor de \(6\).}{}{}}
\End

\Data\ID{9000023709}
\Updated{2022-04-23 05:04:31}
\Level{A}
\SubArea{Radical Equations and Inequalities}
\LANGen{
\Question{
Identify a true statement which concerns the following pair of equations.
\[ 
\begin{aligned}
 \sqrt{
5 - x} & = 2 &\text{(1)}
 \\
\sqrt{x + 5} & = 4 &\text{(2)}
 \end{aligned}
\]}
{The solution of (1) is smaller than the solution of (2).}{The solutions of both equations are prime numbers.}{The solution of (1) is bigger than the solution of (2).}{The solution of (1) equals to the solution of (2).}{}{}}
\LANGpl{
\Question{
Które z poniższych stwierdzeń dotyczących danej pary  równań  jest prawdziwe?
\[ 
\begin{aligned}
 \sqrt{
5 - x} & = 2 &\text{(1)}
 \\
\sqrt{x + 5} & = 4 &\text{(2)}
 \end{aligned}
\]}
{Rozwiązanie (1) jest mniejsze niż rozwiązanie (2).}{Rozwiązanie obydwu równań są liczbami pierwszymi.}{Rozwiązanie  (1) jest większe niż rozwiązanie (2).}{Rozwiązanie (1) jest równe rozwiązaniu (2).}{}{}}
\LANGcs{
\Question{
Jsou dány rovnice.
\[ 
\begin{aligned}
 \sqrt{
5 - x} & = 2 &\text{(1)}
 \\\sqrt{x + 5} & = 4 &\text{(2)}
 \end{aligned}
\]
Které z následujících tvrzení je správné?}
{Číslo, které je řešením rovnice (1), je menší než číslo, které je
řešením rovnice (2).}{Řešením obou rovnic jsou prvočísla.}{Číslo, které je řešením rovnice (1), je větší než číslo, které je
řešením rovnice (2).}{Číslo, které je řešením rovnice (1), je rovno číslu, které je
řešením rovnice (2).}{}{}}
\LANGsk{
\Question{
Sú dané rovnice.
\[ 
\begin{aligned}
 \sqrt{
5 - x} & = 2 &\text{(1)}
 \\
\sqrt{x + 5} & = 4 &\text{(2)}
 \end{aligned}
\]  Ktoré z nasledujúcich tvrdení je správne?}
{Koreň rovnice (1) je menší ako koreň rovnice (2).}{Korene obidvoch rovníc sú prvočísla.}{Koreň rovnice (1) je väčší ako koreň rovnice (2).}{Koreň rovnice (1) sa rovná koreňu rovnice (2).}{}{}}
\LANGes{
\Question{
Dadas las ecuaciones:
\[ 
\begin{aligned}
 \sqrt{
5 - x} & = 2 &\text{(1)}
 \\
\sqrt{x + 5} & = 4 &\text{(2)}
 \end{aligned}
\]
Identifica la proposición lógica que hace referencia a las ecuaciones.}
{La solución de la ecuación (1) es menor que la solución de la ecuación (2).}{Las soluciones de ambas ecuaciones son números primos.}{La solución de la ecuación (1) es mayor que la solución de la ecuación (2).}{La solución de la ecuación (1) equivale a la solución de la ecuación (2).}{}{}}
\End

\Data\ID{9000023710}
\Updated{2020-12-29 09:12:27}
\Level{A}
\SubArea{Radical Equations and Inequalities}
\LANGen{
\Question{
Identify a true statement which concerns the following pair of equations.
\[ 
\begin{aligned}
 \sqrt{
2x + 17} & = 3 &\text{(1)}
 \\
\sqrt{8 - 4x} & = 4 &\text{(2)}
 \end{aligned}
\]}
{The product of the solutions of (1) and (2) is
\(8\).}{The sum of the solutions of (1) and (2) is
\(- 2\).}{The quotient of the solution of (1) divided by the solution of (2) is
\(- 2\).}{The quotient of the solution of (2) divided by the solution of (1) is
\(- 0.5\).}{}{}}
\LANGpl{
\Question{
Które z poniższych stwierdzeń dotyczących podanej pary równań jest prawdziwe? 
\[ 
\begin{aligned}
 \sqrt{
2x + 17} & = 3 &\text{(1)}
 \\
\sqrt{8 - 4x} & = 4 &\text{(2)}
 \end{aligned}
\]}
{Produkt rozwiązań równania (1) i (2) jest równa
\(8\).}{Suma rozwiązań równania (1) i (2) jest równa
\(- 2\).}{Iloraz rozwiązania równania (1) podzielonego przez rozwiązanie równania (2) jest równy
\(- 2\).}{Iloraz rozwiązania równania (2) podzielonego przez rozwiązanie równania (1) jest równy
\(- 0.5\).}{}{}}
\LANGcs{
\Question{
Jsou dány rovnice.
\[ 
\begin{aligned}
 \sqrt{
2x + 17} & = 3 &\text{(1)}
 \\
\sqrt{8 - 4x} & = 4 &\text{(2)}
 \end{aligned}
\]
Které z následujících tvrzení je správné?}
{Součin kořenů uvedených rovnic je roven číslu
\(8\).}{Součet kořenů uvedených rovnic je roven číslu
\(- 2\).}{Podíl kořenu rovnice (1) a kořenu rovnice (2) je roven
\(- 2\).}{Podíl kořenu rovnice (2) a kořenu rovnice (1) je roven
\(- 0{,}5\).}{}{}}
\LANGsk{
\Question{
Sú dané rovnice.
\[ 
\begin{aligned}
 \sqrt{
2x + 17} & = 3 &\text{(1)}
 \\
\sqrt{8 - 4x} & = 4 &\text{(2)}
 \end{aligned}
\]  Ktoré z nasledujúcich tvrdení je správne?}
{Súčin koreňov daných rovníc (1) a (2) sa rovná číslu
\(8\).}{Súčet koreňov daných rovníc (1) a (2) sa rovná číslu
\(- 2\).}{Podiel koreňa rovnice (1) a koreňa rovnice (2) sa rovná číslu
\(- 2\).}{Podiel koreňa rovnice (2) a koreňa rovnice (1) sa rovná číslu
\(- 0.5\).}{}{}}
\LANGes{
\Question{
Dadas las ecuaciones:
\[ 
\begin{aligned}
 \sqrt{
2x + 17} & = 3 &\text{(1)}
 \\
\sqrt{8 - 4x} & = 4 &\text{(2)}
 \end{aligned}
\]
Identifica la proposición lógica que hace referencia a las ecuaciones.}
{El producto de las soluciones de ambas ecuaciones es \(8\).}{La suma de las soluciones de ambas ecuaciones es \(- 2\).}{El cociente de la solución de la ecuación (1) y de la solución de la ecuación (2) es \(- 2\).}{El cociente de la solución de la ecuación (2) y de la solución de la ecuación (1) es \(- 0.5\).}{}{}}
\End

\Data\ID{9000023801}
\Updated{2020-07-15 03:07:34}
\Level{A}
\SubArea{Radical Equations and Inequalities}
\LANGen{
\Question{
Find the sum of the solutions of the following equation. 
\[ 
 \sqrt{x - 2} = \frac{x} 
{3}
\]}
{\(9\)}{\(3\)}{\(6\)}{\(12\)}{}{}}
\LANGpl{
\Question{
Oblicz sumę rozwiązań podanego równania.
\[ 
 \sqrt{x - 2} = \frac{x} 
{3}
\]}
{\(9\)}{\(3\)}{\(6\)}{\(12\)}{}{}}
\LANGcs{
\Question{
Je dána rovnice \(\sqrt{x - 2} = \frac{x} 
{3} \).
Určete součet kořenů rovnice.}
{\(9\)}{\(3\)}{\(6\)}{\(12\)}{}{}}
\LANGsk{
\Question{
Určte súčet koreňov danej rovnice. 
\[ 
 \sqrt{x - 2} = \frac{x} 
{3}
\]}
{\(9\)}{\(3\)}{\(6\)}{\(12\)}{}{}}
\LANGes{
\Question{
Halla la suma de las soluciones de la siguiente ecuación:
\[ 
 \sqrt{x - 2} = \frac{x} 
{3}
\]}
{\(9\)}{\(3\)}{\(6\)}{\(12\)}{}{}}
\End

\Data\ID{9000023802}
\Updated{2020-07-15 03:07:34}
\Level{A}
\SubArea{Radical Equations and Inequalities}
\LANGen{
\Question{
Find the product of the solutions of the following equation. 
\[ 
 \sqrt{3x - 8} = \frac{x} 
{2}
\]}
{\(32\)}{\(4\)}{\(8\)}{\(16\)}{}{}}
\LANGpl{
\Question{
Oblicz iloczyn rozwiązań podanego równania.
\[ 
 \sqrt{3x - 8} = \frac{x} 
{2}
\]}
{\(32\)}{\(4\)}{\(8\)}{\(16\)}{}{}}
\LANGcs{
\Question{
Je dána rovnice \(\sqrt{3x - 8} = \frac{x} 
{2} \).
Určete součin kořenů dané rovnice.}
{\(32\)}{\(4\)}{\(8\)}{\(16\)}{}{}}
\LANGsk{
\Question{
Určte súčin koreňov danej rovnice. 
\[ 
 \sqrt{3x - 8} = \frac{x} 
{2}
\]}
{\(32\)}{\(4\)}{\(8\)}{\(16\)}{}{}}
\LANGes{
\Question{
Halla el producto de las soluciones de la siguiente ecuación:
\[ 
 \sqrt{3x - 8} = \frac{x} 
{2}
\]}
{\(32\)}{\(4\)}{\(8\)}{\(16\)}{}{}}
\End

\Data\ID{9000023803}
\Updated{2020-12-29 09:12:01}
\Level{A}
\SubArea{Radical Equations and Inequalities}
\LANGen{
\Question{
In the following list identify a true statement referring to the solution of the
following equation. 
\[ 
 \sqrt{x + 3} = 3 + x
\]}
{The difference of the bigger and smaller solutions is
\(1\).}{The difference of the bigger and smaller solutions is
\(- 1\).}{The difference of the smaller and the bigger solutions is
\(1\).}{The difference of the smaller and twice the bigger solutions is
\(- 1\).}{}{}}
\LANGpl{
\Question{
Które z poniższych stwierdzeń dotyczących rozwiązania podanego równania jest prawdziwe?
\[ 
 \sqrt{x + 3} = 3 + x
\]}
{Różnica pomiędzy większym i mniejszym rozwiązaniem równania wynosi
\(1\).}{Różnica pomiędzy większym i mniejszym rozwiązaniem równania wynosi
\(- 1\).}{Różnica pomiędzy mniejszym i większym rozwiązaniem równania wynosi
\(1\).}{Różnica pomiędzy mniejszym i większym rozwiązaniem równania wynosi
\(- 1\).}{}{}}
\LANGcs{
\Question{
Je dána rovnice \(\sqrt{x + 3} = 3 + x\).
Které z následujících tvrzení je správné?}
{Rozdíl většího a menšího z kořenů této rovnice je roven číslu
\(1\).}{Rozdíl většího a menšího z kořenů této rovnice je roven číslu
\(- 1\).}{Rozdíl menšího a většího z kořenů této rovnice je roven číslu
\(1\).}{Rozdíl menšího a dvojnásobku většího z kořenů této rovnice je roven
číslu \(- 1\).}{}{}}
\LANGsk{
\Question{
Je daná rovnica. 
\[ 
 \sqrt{x + 3} = 3 + x
\]  Ktoré z nasledujúcich tvrdení je pravdivé?}
{Rozdiel väčšieho a menšieho z koreňov tejto rovnice sa rovná číslu
\(1\).}{Rozdiel väčšieho a menšieho z koreňov tejto rovnice sa rovná číslu
\(- 1\).}{Rozdiel menšieho a väčšieho z koreňov tejto rovnice sa rovná číslu
\(1\).}{Rozdiel menšieho a dvojnásobku väčšieho z koreňov tejto rovnice sa rovná číslu
\(- 1\).}{}{}}
\LANGes{
\Question{
Dada la ecuación:
\[ 
 \sqrt{x + 3} = 3 + x
\]
Identifica la proposición lógica que hace referencia a la ecuación.}
{El resto de la mayor solución y de la menor equivale a \(1\).}{El resto de la mayor solución y de la menor equivale a \(- 1\).}{El resto de la menor solución y de la mayor equivale a \(1\).}{El resto de la menor solución y el doble de la mayor solución equivale a \(- 1\).}{}{}}
\End

\Data\ID{9000023804}
\Updated{2020-12-29 09:12:21}
\Level{A}
\SubArea{Radical Equations and Inequalities}
\LANGen{
\Question{
Identify a true statement which concerns to the following equation.
\[ 
 \sqrt{x + 3} = x - 3
\]}
{The solution is in the interval \((5;8)\).}{The solution is in the interval \([ - 2;2] \).}{The solution is in the interval \([ - 3;1)\).}{The solution is in the interval \([ 3;5)\).}{}{}}
\LANGpl{
\Question{
Które z podanych stwierdzeń dotyczących podanego równania jest prawdziwe?
\[ 
 \sqrt{x + 3} = x - 3
\]}
{Rozwiązanie mieści się w przedziale  \((5;8)\).}{Rozwiązanie mieści się w przedziale \([ - 2;2] \).}{Rozwiązanie mieści się w przedziale \([ - 3;1)\).}{Rozwiązanie mieści się w przedziale \([ 3;5)\).}{}{}}
\LANGcs{
\Question{
Je dána rovnice \(\sqrt{x + 3} = x - 3\).
Které z následujících tvrzení je správné?}
{Řešením této rovnice je číslo z intervalu
\((5;8)\).}{Řešením této rovnice je číslo z intervalu
\(\langle - 2;2\rangle \).}{Řešením této rovnice je číslo z intervalu
\(\langle - 3;1)\).}{Řešením této rovnice je číslo z intervalu
\(\langle 3;5)\).}{}{}}
\LANGsk{
\Question{
Je daná rovnica.
\[ 
 \sqrt{x + 3} = x - 3
\] Ktoré z nasledujúcich tvrdení je správne?}
{Riešením tejto rovnice je číslo z intervalu \((5;8)\).}{Riešením tejto rovnice je číslo z intervalu \([ - 2;2] \).}{Riešením tejto rovnice je číslo z intervalu \([ - 3;1)\).}{Riešením tejto rovnice je číslo z intervalu \([ 3;5)\).}{}{}}
\LANGes{
\Question{
Dada la ecuación:
\[ 
 \sqrt{x + 3} = x - 3
\]
Identifica la proposición lógica que hace referencia a la ecuación.}
{La solución es un número del intervalo \((5;8)\).}{La solución es un número del intervalo \([ - 2;2] \).}{La solución es un número del intervalo \([ - 3;1)\).}{La solución es un número del intervalo \([ 3;5)\).}{}{}}
\End

\Data\ID{9000023805}
\Updated{2022-04-23 05:04:31}
\Level{A}
\SubArea{Radical Equations and Inequalities}
\LANGen{
\Question{
Identify a true statement about the following equation.
\[ 
 \sqrt{6 + x} = -x
\]}
{The solution is in the set \(\left \{x\in \mathbb{R} : -4 < x\leq - 1\right \}\).}{The solution is in the set \(\left \{x\in \mathbb{R} : 1\leq x\leq 5\right \}\).}{The solution is in the set \(\left \{x\in \mathbb{R} : -6\leq x\leq - 3\right \}\).}{The solution is in the set \(\left \{x\in \mathbb{R} : -2 < x < 3\right \}\).}{}{}}
\LANGpl{
\Question{
Które z poniższych stwierdzeń dotyczących podanego równania jest prawdziwe?
\[ 
 \sqrt{6 + x} = -x
\]}
{Rozwiązanie należy do zbioru \(\left \{x\in \mathbb{R} : -4 < x\leq - 1\right \}\).}{Rozwiązanie należy do zbioru \(\left \{x\in \mathbb{R} : 1\leq x\leq 5\right \}\).}{Rozwiązanie należy do zbioru  \(\left \{x\in \mathbb{R} : -6\leq x\leq - 3\right \}\).}{Rozwiązanie należy do zbioru \(\left \{x\in \mathbb{R} : -2 < x < 3\right \}\).}{}{}}
\LANGcs{
\Question{
Je dána rovnice \(\sqrt{6 + x} = -x\).
Které z následujících tvrzení je správné?}
{Řešením této rovnice je číslo z množiny
\(\left \{x\in \mathbb{R} : -4 < x\leq - 1\right \}\).}{Řešením této rovnice je číslo z množiny
\(\left \{x\in \mathbb{R} : 1\leq x\leq 5\right \}\).}{Řešením této rovnice je číslo z množiny
\(\left \{x\in \mathbb{R} : -6\leq x\leq - 3\right \}\).}{Řešením této rovnice je číslo z množiny
\(\left \{x\in \mathbb{R} : -2 < x < 3\right \}\).}{}{}}
\LANGsk{
\Question{
Je daná rovnica.
\[ 
 \sqrt{6 + x} = -x
\] Ktoré tvrdenie je správne?}
{Riešením tejto rovnice je číslo z množiny \(\left \{x\in \mathbb{R} : -4 < x\leq - 1\right \}\).}{Riešením tejto rovnice je číslo z množiny \(\left \{x\in \mathbb{R} : 1\leq x\leq 5\right \}\).}{Riešením tejto rovnice je číslo z množiny \(\left \{x\in \mathbb{R} : -6\leq x\leq - 3\right \}\).}{Riešením tejto rovnice je číslo z množiny \(\left \{x\in \mathbb{R} : -2 < x < 3\right \}\).}{}{}}
\LANGes{
\Question{
Dada la ecuación:
\[ 
 \sqrt{6 + x} = -x
\]
Identifica la proposición lógica que hace referencia a la ecuación.}
{La solución es un número del conjunto de soluciones \(\left \{x\in \mathbb{R} : -4 < x\leq - 1\right \}\).}{La solución es un número del conjunto de soluciones \(\left \{x\in \mathbb{R} : 1\leq x\leq 5\right \}\).}{La solución es un número del conjunto de soluciones \(\left \{x\in \mathbb{R} : -6\leq x\leq - 3\right \}\).}{La solución es un número del conjunto de soluciones \(\left \{x\in \mathbb{R} : -2 < x < 3\right \}\).}{}{}}
\End

\Data\ID{9000023806}
\Updated{2020-12-29 09:12:23}
\Level{A}
\SubArea{Radical Equations and Inequalities}
\LANGen{
\Question{
Identify a true statement which concerns to the following equation.
\[ 
 \sqrt{3x + 4} = x
\]}
{The solution divides \(4\).}{The solution divides \(1\).}{The solution divides \(2\).}{The solution divides \(3\).}{}{}}
\LANGpl{
\Question{
Które z poniższych stwierdzeń dotyczących podanego równania jest prawdziwe?
\[ 
 \sqrt{3x + 4} = x
\]}
{Rozwiązanie jest podzielne przez \(4\).}{Rozwiązanie jest podzielne przez \(1\).}{Rozwiązanie jest podzielne przez \(2\).}{Rozwiązanie jest podzielne przez \(3\).}{}{}}
\LANGcs{
\Question{
Je dána rovnice \(\sqrt{3x + 4} = x\).
Které z následujících tvrzení je správné?}
{Řešením této rovnice je číslo, které je dělitelem čísla
\(4\).}{Řešením této rovnice je číslo, které je dělitelem čísla
\(1\).}{Řešením této rovnice je číslo, které je dělitelem čísla
\(2\).}{Řešením této rovnice je číslo, které je dělitelem čísla
\(3\).}{}{}}
\LANGsk{
\Question{
Je daná rovnica.
\[ 
 \sqrt{3x + 4} = x
\] Ktoré z nasledujúcich tvrdení je pravdivé?}
{Riešením tejto rovnice je číslo, ktoré je deliteľom čísla \(4\).}{Riešením tejto rovnice je číslo, ktoré je deliteľom čísla \(1\).}{Riešením tejto rovnice je číslo, ktoré je deliteľom čísla \(2\).}{Riešením tejto rovnice je číslo, ktoré je deliteľom čísla \(3\).}{}{}}
\LANGes{
\Question{
Dada la ecuación:
\[ 
 \sqrt{3x + 4} = x
\]
Identifica la proposición lógica que hace referencia a la ecuación.}
{La solución es un divisor del número \(4\).}{La solución es un divisor del número \(1\).}{La solución es un divisor del número \(2\).}{La solución es un divisor del número \(3\).}{}{}}
\End

\Data\ID{9000023807}
\Updated{2020-12-29 09:12:47}
\Level{A}
\SubArea{Radical Equations and Inequalities}
\LANGen{
\Question{
Identify a true statement which concerns to the following equation.
\[ 
 \sqrt{x + 3} = \frac{x} 
{2}
\]}
{The solution is a multiple of \(2\).}{The solution is a multiple of \(4\).}{The solution is a multiple of \(8\).}{The solution is a multiple of \(12\).}{}{}}
\LANGpl{
\Question{
Które z poniższych stwierdzeń dotyczących podanego równania jest prawdziwe?
\[ 
 \sqrt{x + 3} = \frac{x} 
{2}
\]}
{Rozwiązanie jest wielokrotnością liczby  \(2\).}{Rozwiązanie jest wielokrotnością liczby \(4\).}{Rozwiązanie jest wielokrotnością liczby \(8\).}{Rozwiązanie jest wielokrotnością liczby  \(12\).}{}{}}
\LANGcs{
\Question{
Je dána rovnice \(\sqrt{x + 3} = \frac{x} 
{2} \).
Které z následujících tvrzení je správné?}
{Řešením této rovnice je číslo, které je násobkem čísla
\(2\).}{Řešením této rovnice je číslo, které je násobkem čísla
\(4\).}{Řešením této rovnice je číslo, které je násobkem čísla
\(8\).}{Řešením této rovnice je číslo, které je násobkem čísla
\(12\).}{}{}}
\LANGsk{
\Question{
Je daná rovnica.
\[ 
 \sqrt{x + 3} = \frac{x} 
{2}
\] Ktoré z nasledujúcich tvrdení je pravdivé?}
{Riešením tejto rovnice je číslo, ktoré je násobkom čísla \(2\).}{Riešením tejto rovnice je číslo, ktoré je násobkom čísla \(4\).}{Riešením tejto rovnice je číslo, ktoré je násobkom čísla \(8\).}{Riešením tejto rovnice je číslo, ktoré je násobkom čísla \(12\).}{}{}}
\LANGes{
\Question{
Dada la ecuación:
\[ 
 \sqrt{x + 3} = \frac{x} 
{2}
\]
Identifica la proposición lógica que hace referencia a la ecuación.}
{La solución es un múltiplo de \(2\).}{La solución es un múltiplo de \(4\).}{La solución es un múltiplo de \(8\).}{La solución es un múltiplo de \(12\).}{}{}}
\End

\Data\ID{9000023808}
\Updated{2020-12-29 09:12:04}
\Level{A}
\SubArea{Radical Equations and Inequalities}
\LANGen{
\Question{
Identify a true statement which concerns to the following equation.
\[ 
 \sqrt{x + 5} = x + 3
\]}
{The solution \(x\) 
satisfies \(|x| = 1\).}{The solution \(x\) 
satisfies \(|x| = 2\).}{The solution \(x\) 
satisfies \(|x| = 3\).}{The solution \(x\) 
satisfies \(|x| = 4\).}{}{}}
\LANGpl{
\Question{
Które z poniższych stwierdzeń dotyczących podanego równania jest prawdziwe?
\[ 
 \sqrt{x + 5} = x + 3
\]}
{Rozwiązaniem jest \(|x| = 1\).}{Rozwiązaniem jest  \(|x| = 2\).}{Rozwiązanie jest \(|x| = 3\).}{Rozwiązaniem jest  \(|x| = 4\).}{}{}}
\LANGcs{
\Question{
Je dána rovnice \(\sqrt{x + 5} = x + 3\).
Které z následujících tvrzení je správné?}
{Pro řešení $x$
platí, že \(|x| = 1\).}{Pro řešení $x$
platí, že \(|x| = 2\).}{Pro řešení $x$
platí, že \(|x| = 3\).}{Pro řešení $x$
platí, že \(|x| = 4\).}{}{}}
\LANGsk{
\Question{
Je daná rovnica.
\[ 
 \sqrt{x + 5} = x + 3
\] Ktoré z nasledujúcich tvrdení je pravdivé?}
{Riešením tejto rovnice je číslo \(x\), pre ktoré platí
 \(|x| = 1\).}{Riešením tejto rovnice je číslo \(x\),
pre ktoré platí \(|x| = 2\).}{Riešením tejto rovnice je číslo \(x\), pre ktoré platí
 \(|x| = 3\).}{Riešením tejto rovnice je číslo \(x\), pre ktoré platí
 \(|x| = 4\).}{}{}}
\LANGes{
\Question{
Dada la ecuación:
\[ 
 \sqrt{x + 5} = x + 3
\]
Identifica la proposición lógica que hace referencia a la ecuación.}
{La solución \(x\) 
satisface \(|x| = 1\).}{La solución \(x\) 
satisface \(|x| = 2\).}{La solución \(x\) 
satisface \(|x| = 3\).}{La solución \(x\) 
satisface \(|x| = 4\).}{}{}}
\End

\Data\ID{9000023809}
\Updated{2020-12-29 09:12:26}
\Level{A}
\SubArea{Radical Equations and Inequalities}
\LANGen{
\Question{
Identify a true statement which concerns to the following equation.
\[ 
 \sqrt{16 - 5x} = 2 - x
\]}
{The solution \(x\) 
satisfies \(|x| > 3\).}{The solution \(x\) 
satisfies \(|x| < 3\).}{The solution \(x\) 
satisfies \(|x + 1| < 3\).}{The solution \(x\) 
satisfies \(|x + 1| > 3\).}{}{}}
\LANGpl{
\Question{
Które z poniższych stwierdzeń dotyczących podanego równania jest prawdziwe?
\[ 
 \sqrt{16 - 5x} = 2 - x
\]}
{Rozwiązaniem równania jest \(|x| > 3\).}{Rozwiązaniem równania jest \(|x| < 3\).}{Rozwiązaniem równania jest  \(|x + 1| < 3\).}{Rozwiązaniem równania jest  \(|x + 1| > 3\).}{}{}}
\LANGcs{
\Question{
Je dána rovnice \(\sqrt{16 - 5x} = 2 - x\).
Které z následujících tvrzení je správné?}
{Pro řešení \(x\) platí, že \(|x| > 3\).}{Pro řešení \(x\) platí, že \(|x| < 3\).}{Pro řešení \(x\) platí, že \(|x + 1| < 3\).}{Pro řešení \(x\) platí, že \(|x + 1| > 3\).}{}{}}
\LANGsk{
\Question{
Je daná rovnica.
\[ 
 \sqrt{16 - 5x} = 2 - x
\] Ktoré z nasledujúcich tvrdení je správne?}
{Riešením tejto rovnice je číslo \(x\), pre ktoré platí
 \(|x| > 3\).}{Riešením tejto rovnice je číslo \(x\), pre ktoré platí 
 \(|x| < 3\).}{Riešením tejto rovnice je číslo \(x\),
pre ktoré platí \(|x + 1| < 3\).}{Riešením tejto rovnice je číslo \(x\),
pre ktoré platí \(|x + 1| > 3\).}{}{}}
\LANGes{
\Question{
Dada la ecuación:
\[ 
 \sqrt{16 - 5x} = 2 - x
\]
Identifica la proposición lógica que hace referencia a la ecuación.}
{La solución \(x\) 
satisface \(|x| > 3\).}{La solución \(x\) 
satisface \(|x| < 3\).}{La solución \(x\) 
satisface \(|x + 1| < 3\).}{La solución \(x\) 
satisface \(|x + 1| > 3\).}{}{}}
\End

\Data\ID{9000023810}
\Updated{2020-07-15 03:07:34}
\Level{A}
\SubArea{Radical Equations and Inequalities}
\LANGen{
\Question{
Denote by \(x_{1}\) 
the solution of the equation 
\[ 
 \sqrt{6 - 2x} = -x - 1
\]
 and by \(x_{2}\) 
the solution of the equation 
\[ 
 \sqrt{2x + 6} = 9 - x.
\]
 Identify a correct statement about \(x_{1}\) 
and \(x_{2}\).}
{\(|x_{1}| = |x_{2}|\)}{\(|x_{1}| < |x_{2}|\)}{\(|x_{1}| > |x_{2}|\)}{\(5|x_{1}| = |x_{2}|\)}{}{}}
\LANGpl{
\Question{
Oznaczając rozwiązanie  równania \[ 
 \sqrt{6 - 2x} = -x - 1
\]
 symbolem  \(x_{1}\) a rozwiązanie innego równania symbolem  \(x_{2}\) 
\[ 
 \sqrt{2x + 6} = 9 - x.
\]
Wybierz poprawną odpowiedź dotyczącą rozwiązań  \(x_{1}\) 
i \(x_{2}\).}
{\(|x_{1}| = |x_{2}|\)}{\(|x_{1}| < |x_{2}|\)}{\(|x_{1}| > |x_{2}|\)}{\(5|x_{1}| = |x_{2}|\)}{}{}}
\LANGcs{
\Question{
Jsou dány rovnice
\[ \begin{aligned}
 \sqrt{
6 - 2x} & = -x - 1, &\text{(1)}
 \\\sqrt{2x + 6} & = 9 - x. &\text{(2)}
 \\\end{aligned}\]
Řešení rovnice (1) označme \(x_{1}\) a
řešení rovnice (2) označme \(x_{2}\).
Které z následujících tvrzení je správné?}
{\(|x_{1}| = |x_{2}|\)}{\(|x_{1}| < |x_{2}|\)}{\(|x_{1}| > |x_{2}|\)}{\(5|x_{1}| = |x_{2}|\)}{}{}}
\LANGsk{
\Question{
Označme \(x_{1}\) 
riešenie rovnice 
\[ 
 \sqrt{6 - 2x} = -x - 1
\]
 a  \(x_{2}\) 
riešenie rovnice 
\[ 
 \sqrt{2x + 6} = 9 - x.
\]
 Ktoré z nasledujúcich tvrdení o  \(x_{1}\) 
a \(x_{2}\) je správne?}
{\(|x_{1}| = |x_{2}|\)}{\(|x_{1}| < |x_{2}|\)}{\(|x_{1}| > |x_{2}|\)}{\(5|x_{1}| = |x_{2}|\)}{}{}}
\LANGes{
\Question{
Denota por \(x_{1}\) 
la solución de la ecuación 
\[ 
 \sqrt{6 - 2x} = -x - 1
\]
y por \(x_{2}\) 
la solución de la ecuación 
\[ 
 \sqrt{2x + 6} = 9 - x.
\]
Identifica la proposición lógica sobre \(x_{1}\) 
y \(x_{2}\).}
{\(|x_{1}| = |x_{2}|\)}{\(|x_{1}| < |x_{2}|\)}{\(|x_{1}| > |x_{2}|\)}{\(5|x_{1}| = |x_{2}|\)}{}{}}
\End

\Data\ID{9000024802}
\Updated{2020-12-29 09:12:18}
\Level{A}
\SubArea{Radical Equations and Inequalities}
\LANGen{
\Question{
Consider the equation 
\[ 
 \sqrt{x^{2 } - 2x + 1} = x + 2
\]
 and the equation which arises from this equation by squaring both sides of the
equation, i.e. the equation 
\[ 
 \left (\sqrt{x^{2 } - 2x + 1}\right )^{2} = (x + 2)^{2}.
\]
 Identify a true statement.}
{Both equations are equivalent only if
\(x\geq - 2\).}{Both equations are equivalent.}{Both equations are equivalent only if
\(x\leq - 2\).}{None of the above.}{}{}}
\LANGpl{
\Question{
Rozważ równanie
\[ 
 \sqrt{x^{2 } - 2x + 1} = x + 2
\]
i równanie, które powstało poprzez podniesienie do kwadratu  obu stron podanego równania 
\[ 
 \left (\sqrt{x^{2 } - 2x + 1}\right )^{2} = (x + 2)^{2}.
\]
Które zdanie jest prawdziwe?}
{Obydwa równania są równoważne tylko wtedy, gdy
\(x\geq - 2\).}{Obydwa równania są równoważne.}{Obydwa równania są równoważne tylko wtedy, gdy 
\(x\leq - 2\).}{Żadne z powyższych stwierdzeń nie jest prawdą.}{}{}}
\LANGcs{
\Question{
Uvažujme rovnici 
\[ 
 \sqrt{x^{2 } - 2x + 1} = x + 2
\]
 a rovnici, která z této rovnice vznikne umocněním obou stran rovnice na
druhou, tj. rovnici 
\[ 
 \left (\sqrt{x^{2 } - 2x + 1}\right )^{2} = (x + 2)^{2}.
\]
 Označte správné tvrzení.}
{Obě rovnice jsou ekvivalentní pouze pro
\(x\geq - 2\).}{Obě rovnice jsou ekvivalentní.}{Obě rovnice jsou ekvivalentní pouze pro
\(x\leq - 2\).}{Žádná z výše uvedených odpovědí není správná.}{}{}}
\LANGsk{
\Question{
Uvažujme o rovnici
\[ 
 \sqrt{x^{2 } - 2x + 1} = x + 2
\]
 a o rovnici, ktorá z tejto rovnice vznikne umocnením obidvoch strán rovnice na druhú, tj. o rovnici 
\[ 
 \left (\sqrt{x^{2 } - 2x + 1}\right )^{2} = (x + 2)^{2}.
\]
 Označte správne tvrdenie.}
{Obidve rovnice sú ekvivalentné len pre
\(x\geq - 2\).}{Obidve rovnice sú ekvivalentné.}{Obidve rovnice sú ekvivalentné len pre
\(x\leq - 2\).}{Žiadna z vyššie uvedených odpovedí nie je správna.}{}{}}
\LANGes{
\Question{
Consideramos la ecuación
\[ 
 \sqrt{x^{2 } - 2x + 1} = x + 2
\]
 y la ecuación que surge de esta ecuación al elevar al cuadrado ambos lados de la ecuación, es decir, la ecuación: 
\[ 
 \left (\sqrt{x^{2 } - 2x + 1}\right )^{2} = (x + 2)^{2}.
\]
Identifica la proposición lógica que hace referencia a estas ecuaciones}
{Ambas ecuaciones son equivalentes solo si \(x\geq - 2\).}{Ambas ecuaciones son equivalentes.}{Ambas ecuaciones son equivalentes solo si \(x\leq - 2\).}{Ninguna de las respuestas es correcta.}{}{}}
\End

\Data\ID{9000024803}
\Updated{2020-12-29 09:12:09}
\Level{A}
\SubArea{Radical Equations and Inequalities}
\LANGen{
\Question{
Removing radical in an equation by squaring both sides may enrich the set of
solutions of this equation and checking the solutions of the new equation in the
original equation may be necessary. Identify a correct conclusion in the particular
case of the following equation. 
\[ 
 -\sqrt{x^{2 } - 2x + 1} = x
\]}
{If we look for the solution in the set
\(\mathbb{R}^{-}\), then
squaring both sides of the equation gives an equivalent equation. The checking of the
solution is not necessary.}{If we look for the solution in the set
\(\mathbb{R}^{+}\), then
squaring both sides of the equation gives an equivalent equation. The checking of the
solution is not necessary.}{If we look for the solution in the set
\(\mathbb{R}\), then
squaring both sides of the equation gives an equivalent equation. The checking of the
solution is not necessary.}{None of the above.}{}{}}
\LANGpl{
\Question{
Usunięcie pierwiastka z równania poprzez podniesienie obydwu stron tego równania do kwadratu może wzbogacić zbiór rozwiązań tego równania. Konieczne będzie jednak sprawdzenie rozwiązań powstałego równania w pierwotnej postaci tego równania.  Określ, który z poniższych wniosków jest zgodny z prawdą dla podanego równania. \[ 
 -\sqrt{x^{2 } - 2x + 1} = x
\]}
{Jeśli szukamy rozwiązań równania w zbiorze 
\(\mathbb{R}^{-}\), wtedy podniesienie obydwu stron równania do kwadratu daje nam równoważne równanie. Sprawdzanie rozwiązań w tym przypadku nie jest konieczne.}{Jeśli szukamy rozwiązań równania w zbiorze 
\(\mathbb{R}^{+}\), wtedy podniesienie obydwu stron równania do kwadratu daje nam równoważne równanie. Sprawdzanie rozwiązań w tym przypadku nie jest konieczne.}{Jeśli szukamy rozwiązań równania w zbiorze 
\(\mathbb{R}\), wtedy podniesienie obydwu stron równania do kwadratu daje nam równoważne równanie. Sprawdzanie rozwiązań w tym przypadku nie jest konieczne.}{Żadne z powyższych stwierdzeń nie jest prawdziwe.}{}{}}
\LANGcs{
\Question{
Odstranění odmocnin v rovnici umocněním obou stran rovnice na druhou
může rozšířit množinu řešení. Pro kořeny nové rovnice může být
nutné provést zkoušku, zda jsou i kořeny rovnice původní. Rozhodněte o
nutnosti provedení zkoušky v závislosti na definičním oboru při řešení
rovnice.
\[ 
 -\sqrt{x^{2 } - 2x + 1} = x
\]}
{Řešíme-li v \(\mathbb{R}^{-}\),
umocněním obou stran rovnice obdržíme ekvivalentní rovnici a zkouška
není nutnou součástí řešení.}{Řešíme-li v \(\mathbb{R}^{+}\),
umocněním obou stran rovnice obdržíme ekvivalentní rovnici a zkouška
není nutnou součástí řešení.}{Řešíme-li v \(\mathbb{R}\),
umocněním obou stran rovnice obdržíme ekvivalentní rovnici a zkouška
není nutnou součástí řešení.}{Ani jedna z výše uvedených odpovědí není správná.}{}{}}
\LANGsk{
\Question{
Odstránenie odmocnín v rovnici umocnením obidvoch strán rovnice na druhú môže rozšíriť množinu riešení. Pre korene novej rovnice môže byť nutné urobiť skúšku, či sú aj koreňmi pôvodnej rovnice. Rozhodnite o nutnosti prevedenia skúšky v závislosti od definičného oboru pri riešení rovnice. 
\[ 
 -\sqrt{x^{2 } - 2x + 1} = x
\]}
{Ak riešime v
\(\mathbb{R}^{-}\), potom umocnením obidvoch strán rovnice dostaneme ekvivalentnú rovnicu a skúška nie je nutnou súčasťou riešenia.}{Ak riešime v
\(\mathbb{R}^{+}\), potom
umocnením obidvoch strán rovnice dostaneme ekvivalentnú rovnicu a skúška nie je nutnou súčasťou riešenia.}{Ak riešime v
\(\mathbb{R}\), potom
umocnením obidvoch strán rovnice dostaneme ekvivalentnú rovnicu a skúška nie je nutnou súčasťou riešenia.}{Ani jedna z vyššie uvedených odpovedí nie je správna.}{}{}}
\LANGes{
\Question{
Eliminar el radical en una ecuación, al elevar al cuadrado ambos lados, puede ampliar el conjunto de soluciones de esta ecuación y puede ser necesario verificar las soluciones de la nueva ecuación en la ecuación original. Identifica la conclusión correcta en el caso particular de la siguiente ecuación.
\[ 
 -\sqrt{x^{2 } - 2x + 1} = x
\]}
{Si buscamos la solución en el conjunto
\(\mathbb{R}^{-}\), entonces la cuadratura de ambos lados de la ecuación da una ecuación equivalente. La comprobación de la solución no es necesaria.}{Si buscamos la solución en el conjunto
\(\mathbb{R}^{+}\), entonces la cuadratura de ambos lados de la ecuación da una ecuación equivalente. La comprobación de la solución no es necesaria.}{Si buscamos la solución en el conjunto
\(\mathbb{R}\), entonces la cuadratura de ambos lados de la ecuación da una ecuación equivalente. La comprobación de la solución no es necesaria.}{Ninguna respuesta es correcta.}{}{}}
\End

\Data\ID{9000033702}
\Updated{2020-07-15 03:07:34}
\Level{A}
\SubArea{Radical Equations and Inequalities}
\LANGen{
\Question{
Find the domain of the following expression. 
\[ 
 \sqrt{-x^{2 } + 7x - 12} -\frac{1} 
{x}
\]}
{\([ 3;4] \)}{\(\mathbb{R}\setminus \left \{0\right \}\)}{\(\mathbb{R}\setminus \left \{0;3;4\right \}\)}{\(\left (3;4\right )\)}{\(\left (-\infty ;3\right )\cup \left (4;\infty \right )\)}{\(\left (-\infty ;3] \cup [ 4;\infty \right )\)}}
\LANGpl{
\Question{
Znajdź dziedzinę podanego wyrażenia.
\[ 
 \sqrt{-x^{2 } + 7x - 12} -\frac{1} 
{x}
\]}
{\([ 3;4] \)}{\(\mathbb{R}\setminus \left \{0\right \}\)}{\(\mathbb{R}\setminus \left \{0;3;4\right \}\)}{\(\left (3;4\right )\)}{\(\left (-\infty ;3\right )\cup \left (4;\infty \right )\)}{\(\left (-\infty ;3] \cup [ 4;\infty \right )\)}}
\LANGcs{
\Question{
Výraz \(\sqrt{-x^{2 } + 7x - 12} -\frac{1} 
{x}\) 
má definiční obor.}
{\(\langle 3;4\rangle \)}{\(\mathbb{R}\setminus \left \{0\right \}\)}{\(\mathbb{R}\setminus \left \{0;3;4\right \}\)}{\(\left (3;4\right )\)}{\(\left (-\infty ;3\right )\cup \left (4;\infty \right )\)}{\(\left (-\infty ;3\rangle \cup \langle 4;\infty \right )\)}}
\LANGsk{
\Question{
Určte definičný obor nasledujúceho výrazu.
\[ 
 \sqrt{-x^{2 } + 7x - 12} -\frac{1} 
{x}
\]}
{\([ 3;4] \)}{\(\mathbb{R}\setminus \left \{0\right \}\)}{\(\mathbb{R}\setminus \left \{0;3;4\right \}\)}{\(\left (3;4\right )\)}{\(\left (-\infty ;3\right )\cup \left (4;\infty \right )\)}{\(\left (-\infty ;3] \cup [ 4;\infty \right )\)}}
\LANGes{
\Question{
Halla el dominio de la siguiente expresión. 
\[ 
 \sqrt{-x^{2 } + 7x - 12} -\frac{1} 
{x}
\]}
{\([ 3;4] \)}{\(\mathbb{R}\setminus \left \{0\right \}\)}{\(\mathbb{R}\setminus \left \{0;3;4\right \}\)}{\(\left (3;4\right )\)}{\(\left (-\infty ;3\right )\cup \left (4;\infty \right )\)}{\(\left (-\infty ;3] \cup [ 4;\infty \right )\)}}
\End

\Data\ID{9000034901}
\Updated{2022-04-23 05:04:31}
\Level{A}
\SubArea{Radical Equations and Inequalities}
\LANGen{
\Question{
Find the domain of the following expression. 
\[ 
 \sqrt{\left (2x - 3 \right ) \left (3x + 1 \right )}
\]}
{\(\left (-\infty ;-\frac{1}
{3}\right ] \cup \left [ \frac{3}
{2};\infty \right )\)}{\(\left [ -\frac{1}
{3}; \frac{3} 
{2}\right ] \)}{\(\left (-\frac{1}
{3}; \frac{3} 
{2}\right )\)}{\(\left (-\infty ;-\frac{1}
{3}\right )\cup \left (\frac{3}
{2};\infty \right )\)}{}{}}
\LANGpl{
\Question{
Znajdź dziedzinę podanego wyrażenia.
\[ 
 \sqrt{\left (2x - 3 \right ) \left (3x + 1 \right )}
\]}
{\(\left (-\infty ;-\frac{1}
{3}\right ] \cup \left [ \frac{3}
{2};\infty \right )\)}{\(\left [ -\frac{1}
{3}; \frac{3} 
{2}\right ] \)}{\(\left (-\frac{1}
{3}; \frac{3} 
{2}\right )\)}{\(\left (-\infty ;-\frac{1}
{3}\right )\cup \left (\frac{3}
{2};\infty \right )\)}{}{}}
\LANGcs{
\Question{
Množina všech \(x\in \mathbb{R}\),
pro která je výraz \(\sqrt{\left (2x - 3 \right ) \left (3x + 1 \right )}\) 
definován, je:}
{\(\left (-\infty ;-\frac{1}
{3}\right \rangle \cup \left \langle \frac{3}
{2};\infty \right )\)}{\(\left \langle -\frac{1}
{3}; \frac{3} 
{2}\right \rangle \)}{\(\left (-\frac{1}
{3}; \frac{3} 
{2}\right )\)}{\(\left (-\infty ;-\frac{1}
{3}\right )\cup \left (\frac{3}
{2};\infty \right )\)}{}{}}
\LANGsk{
\Question{
Nájdite definičný obor daného výrazu. 
\[ 
 \sqrt{\left (2x - 3 \right ) \left (3x + 1 \right )}
\]}
{\(\left (-\infty ;-\frac{1}
{3}\right ] \cup \left [ \frac{3}
{2};\infty \right )\)}{\(\left [ -\frac{1}
{3}; \frac{3} 
{2}\right ] \)}{\(\left (-\frac{1}
{3}; \frac{3} 
{2}\right )\)}{\(\left (-\infty ;-\frac{1}
{3}\right )\cup \left (\frac{3}
{2};\infty \right )\)}{}{}}
\LANGes{
\Question{
Halla el dominio de la siguiente expresión.
\[ 
 \sqrt{\left (2x - 3 \right ) \left (3x + 1 \right )}
\]}
{\(\left (-\infty ;-\frac{1}
{3}\right ] \cup \left [ \frac{3}
{2};\infty \right )\)}{\(\left [ -\frac{1}
{3}; \frac{3} 
{2}\right ] \)}{\(\left (-\frac{1}
{3}; \frac{3} 
{2}\right )\)}{\(\left (-\infty ;-\frac{1}
{3}\right )\cup \left (\frac{3}
{2};\infty \right )\)}{}{}}
\End

\Data\ID{9000034903}
\Updated{2020-07-15 03:07:34}
\Level{A}
\SubArea{Radical Equations and Inequalities}
\LANGen{
\Question{
Find all \(x\in \mathbb{R}\) 
for which the following expression is undefined. 
\[ 
 \sqrt{\left (3x + 4 \right ) \left (\frac{1}
{5} - x\right )}
\]}
{\(\left (-\infty ;-\frac{4}
{3}\right )\cup \left (\frac{1}
{5};\infty \right )\)}{\(\left [ -\frac{4}
{3}; \frac{1} 
{5}\right ] \)}{\(\left (-\infty ;-\frac{4}
{3}\right ] \cup \left [ \frac{1}
{5};\infty \right )\)}{\(\left (-\frac{4}
{3}; \frac{1} 
{5}\right )\)}{}{}}
\LANGpl{
\Question{
Znajdź wszystkie wartości  \(x\in \mathbb{R}\), dla których podane wyrażenie jest nieokreślone. 
\[ 
 \sqrt{\left (3x + 4 \right ) \left (\frac{1}
{5} - x\right )}
\]}
{\(\left (-\infty ;-\frac{4}
{3}\right )\cup \left (\frac{1}
{5};\infty \right )\)}{\(\left [ -\frac{4}
{3}; \frac{1} 
{5}\right ] \)}{\(\left (-\infty ;-\frac{4}
{3}\right ] \cup \left [ \frac{1}
{5};\infty \right )\)}{\(\left (-\frac{4}
{3}; \frac{1} 
{5}\right )\)}{}{}}
\LANGcs{
\Question{
Množina všech \(x\in \mathbb{R}\), pro
která není výraz \(\sqrt{\left (3x + 4 \right ) \left (\frac{1}
{5} - x\right )}\) 
definován, je:}
{\(\left (-\infty ;-\frac{4}
{3}\right )\cup \left (\frac{1}
{5};\infty \right )\)}{\(\left \langle -\frac{4}
{3}; \frac{1} 
{5}\right \rangle \)}{\(\left (-\infty ;-\frac{4}
{3}\right \rangle \cup \left \langle \frac{1}
{5};\infty \right )\)}{\(\left (-\frac{4}
{3}; \frac{1} 
{5}\right )\)}{}{}}
\LANGsk{
\Question{
Nájdite všetky \(x\in \mathbb{R}\),
pre ktoré nasledujúci výraz nie je definovaný. 
\[ 
 \sqrt{\left (3x + 4 \right ) \left (\frac{1}
{5} - x\right )}
\]}
{\(\left (-\infty ;-\frac{4}
{3}\right )\cup \left (\frac{1}
{5};\infty \right )\)}{\(\left [ -\frac{4}
{3}; \frac{1} 
{5}\right ] \)}{\(\left (-\infty ;-\frac{4}
{3}\right ] \cup \left [ \frac{1}
{5};\infty \right )\)}{\(\left (-\frac{4}
{3}; \frac{1} 
{5}\right )\)}{}{}}
\LANGes{
\Question{
Halla el dominio de la siguiente expresión.
\[ 
 \sqrt{\left (3x + 4 \right ) \left (\frac{1}
{5} - x\right )}
\]}
{\(\left (-\infty ;-\frac{4}
{3}\right )\cup \left (\frac{1}
{5};\infty \right )\)}{\(\left [ -\frac{4}
{3}; \frac{1} 
{5}\right ] \)}{\(\left (-\infty ;-\frac{4}
{3}\right ] \cup \left [ \frac{1}
{5};\infty \right )\)}{\(\left (-\frac{4}
{3}; \frac{1} 
{5}\right )\)}{}{}}
\End

\Data\ID{1003020701}
\Updated{2022-08-25 10:08:13}
\Level{B}
\SubArea{Radical Equations and Inequalities}
\LANGen{
\Question{
Find the solution of the inequality.
\[
\sqrt{-2x^2+x+5}\geq5
\]}
{\(x\in\emptyset\)}{\(x\in[-1;5]\)}{\(x\in\left(-\frac12;\frac52\right)\)}{\(x\in\mathbb{R}\)}{}{}}
\LANGpl{
\Question{
Znajdź rozwiązanie nierówności.
\[
\sqrt{-2x^2+x+5}\geq5
\]}
{\(x\in\emptyset\)}{\(x\in\langle-1;5\rangle\)}{\(x\in\left(-\frac12;\frac52\right)\)}{\(x\in\mathbb{R}\)}{}{}}
\LANGcs{
\Question{
Určete řešení nerovnice.
\[
\sqrt{-2x^2+x+5}\geq5
\]}
{\(x\in\emptyset\)}{\(x\in\langle-1;5\rangle\)}{\(x\in\left(-\frac12;\frac52\right)\)}{\(x\in\mathbb{R}\)}{}{}}
\LANGsk{
\Question{
Vyriešte nerovnicu.
\[
\sqrt{-2x^2+x+5}\geq5
\]}
{\(x\in\emptyset\)}{\(x\in\langle-1;5\rangle\)}{\(x\in\left(-\frac12;\frac52\right)\)}{\(x\in\mathbb{R}\)}{}{}}
\LANGes{
\Question{
Halla la solución de la siguiente inecuación.
\[
\sqrt{-2x^2+x+5}\geq5
\]}
{\(x\in\emptyset\)}{\(x\in[-1;5]\)}{\(x\in\left(-\frac12;\frac52\right)\)}{\(x\in\mathbb{R}\)}{}{}}
\End

\Data\ID{1003020702}
\Updated{2022-08-25 10:08:13}
\Level{B}
\SubArea{Radical Equations and Inequalities}
\LANGen{
\Question{
Find the solution of the inequality.
\[
\sqrt{3x^2-x+7}\geq2
\]}
{\(x\in\mathbb{R}\)}{\(x\in\emptyset\)}{\(x\in\left(-\frac13;\frac73\right)\)}{\(x\in\mathbb{R}\setminus\{2\}\)}{}{}}
\LANGpl{
\Question{
Znajdź rozwiązanie nierówności.
\[
\sqrt{3x^2-x+7}\geq2
\]}
{\(x\in\mathbb{R}\)}{\(x\in\emptyset\)}{\(x\in\left(-\frac13;\frac73\right)\)}{\(x\in\mathbb{R}\setminus\{2\}\)}{}{}}
\LANGcs{
\Question{
Určete řešení nerovnice.
\[
\sqrt{3x^2-x+7}\geq2
\]}
{\(x\in\mathbb{R}\)}{\(x\in\emptyset\)}{\(x\in\left(-\frac13;\frac73\right)\)}{\(x\in\mathbb{R}\setminus\{2\}\)}{}{}}
\LANGsk{
\Question{
Vyriešte nerovnicu.
\[
\sqrt{3x^2-x+7}\geq2
\]}
{\(x\in\mathbb{R}\)}{\(x\in\emptyset\)}{\(x\in\left(-\frac13;\frac73\right)\)}{\(x\in\mathbb{R}\setminus\{2\}\)}{}{}}
\LANGes{
\Question{
Halla la solución de la siguiente inecuación.
\[
\sqrt{3x^2-x+7}\geq2
\]}
{\(x\in\mathbb{R}\)}{\(x\in\emptyset\)}{\(x\in\left(-\frac13;\frac73\right)\)}{\(x\in\mathbb{R}\setminus\{2\}\)}{}{}}
\End

\Data\ID{2010007705}
\Updated{2022-02-15 07:02:37}
\Level{B}
\SubArea{Radical Equations and Inequalities}
\LANGen{
\Question{
Find the solution of the inequality.
\[
\sqrt{-x^2+x+2}\geq 4
\]}
{\(x\in\emptyset\)}{\(x\in [ -3;4]\)}{\(x\in\mathbb{R}\)}{\(x\in [ -1;2]\)}{}{}}
\LANGpl{
\Question{
Rozwiąż nierówność.
\[
\sqrt{-x^2+x+2}\geq 4
\]}
{\(x\in\emptyset\)}{\(x\in \langle -3;4\rangle\)}{\(x\in\mathbb{R}\)}{\(x\in \langle -1;2\rangle\)}{}{}}
\LANGcs{
\Question{
Co je řešením nerovnice?
\[
\sqrt{-x^2+x+2}\geq 4
\]}
{\(x\in\emptyset\)}{\(x\in \langle -3;4\rangle\)}{\(x\in\mathbb{R}\)}{\(x\in \langle -1;2\rangle\)}{}{}}
\LANGsk{
\Question{
Nájdite riešenie nerovnice.
\[
\sqrt{-x^2+x+2}\geq 4
\]}
{\(x\in\emptyset\)}{\(x\in \langle -3;4\rangle\)}{\(x\in\mathbb{R}\)}{\(x\in \langle -1;2\rangle\)}{}{}}
\LANGes{
\Question{
Halla la solución de la siguiente inecuación.
\[
\sqrt{-x^2+x+2}\geq 4
\]}
{\(x\in\emptyset\)}{\(x\in [ -3;4]\)}{\(x\in\mathbb{R}\)}{\(x\in [ -1;2]\)}{}{}}
\End

\Data\ID{9000021807}
\Updated{2020-07-15 03:07:34}
\Level{B}
\SubArea{Radical Equations and Inequalities}
\LANGen{
\Question{
Solve the following inequality. 
\[ 
 \sqrt{\frac{x^{5 } x^{-2 } }
{x^{6}x^{-3}}}\geq 1
\]}
{\(x\in \mathbb{R}\setminus \{0\}\)}{\(x\in \mathbb{R}\)}{\(x\in [ 1;\infty )\)}{\(x\in \emptyset \)}{}{}}
\LANGpl{
\Question{
Rozwiąż podaną nierówność.
\[ 
 \sqrt{\frac{x^{5 } x^{-2 } }
{x^{6}x^{-3}}}\geq 1
\]}
{\(x\in \mathbb{R}\setminus \{0\}\)}{\(x\in \mathbb{R}\)}{\(x\in [ 1;\infty )\)}{\(x\in \emptyset \)}{}{}}
\LANGcs{
\Question{
Množina všech řešení nerovnice
\(\sqrt{
\frac{x^{5}x^{-2}} 
{x^{6}x^{-3}}} \geq 1\) 
je:}
{\(\mathbb{R}\setminus \{0\}\)}{\(\mathbb{R}\)}{\(\langle 1;\infty )\)}{\(\emptyset \)}{}{}}
\LANGsk{
\Question{
Vyriešte danú nerovnicu. 
\[ 
 \sqrt{\frac{x^{5 } x^{-2 } }
{x^{6}x^{-3}}}\geq 1
\]}
{\(x\in \mathbb{R}\setminus \{0\}\)}{\(x\in \mathbb{R}\)}{\(x\in [ 1;\infty )\)}{\(x\in \emptyset \)}{}{}}
\LANGes{
\Question{
Resuelve la siguiente inecuación: 
\[ 
 \sqrt{\frac{x^{5 } x^{-2 } }
{x^{6}x^{-3}}}\geq 1
\]}
{\(x\in \mathbb{R}\setminus \{0\}\)}{\(x\in \mathbb{R}\)}{\(x\in [ 1;\infty )\)}{\(x\in \emptyset \)}{}{}}
\End

\Data\ID{9000024801}
\Updated{2020-07-15 03:07:34}
\Level{B}
\SubArea{Radical Equations and Inequalities}
\LANGen{
\Question{
In the following list identify an inequality which does not have a solution.}
{\(\sqrt{2x - 3} < -6\)}{\(\sqrt{x^{2 } - 3x} > 5\)}{\(\sqrt{1 + x^{2}} > -10\)}{\(\sqrt{2x^{2}} < 4\)}{}{}}
\LANGpl{
\Question{
Z podanych odpowiedzi wybierz nierówność, która nie ma rozwiązania.}
{\(\sqrt{2x - 3} < -6\)}{\(\sqrt{x^{2 } - 3x} > 5\)}{\(\sqrt{1 + x^{2}} > -10\)}{\(\sqrt{2x^{2}} < 4\)}{}{}}
\LANGcs{
\Question{
Která z následujících nerovnic nemá řešení?}
{\(\sqrt{2x - 3} < -6\)}{\(\sqrt{x^{2 } - 3x} > 5\)}{\(\sqrt{1 + x^{2}} > -10\)}{\(\sqrt{2x^{2}} < 4\)}{}{}}
\LANGsk{
\Question{
Ktorá z nasledujúcich nerovníc nemá riešenie?}
{\(\sqrt{2x - 3} < -6\)}{\(\sqrt{x^{2 } - 3x} > 5\)}{\(\sqrt{1 + x^{2}} > -10\)}{\(\sqrt{2x^{2}} < 4\)}{}{}}
\LANGes{
\Question{
¿Cuál de las siguientes inecuaciones no tiene solución?}
{\(\sqrt{2x - 3} < -6\)}{\(\sqrt{x^{2 } - 3x} > 5\)}{\(\sqrt{1 + x^{2}} > -10\)}{\(\sqrt{2x^{2}} < 4\)}{}{}}
\End

\Data\ID{9000024804}
\Updated{2020-07-15 03:07:34}
\Level{B}
\SubArea{Radical Equations and Inequalities}
\LANGen{
\Question{
How many solutions does the inequality 
\[ 
 \sqrt{x + 17} > x - 3
\]
 have in the set \(\mathbb{N}\)?}
{Seven solutions in \(\mathbb{N}\).}{No solution in \(\mathbb{N}\).}{Five solutions in \(\mathbb{N}\).}{More than seven solutions in \(\mathbb{N}\).}{}{}}
\LANGpl{
\Question{
Ile rozwiązań nierówności
\[ 
 \sqrt{x + 17} > x - 3
\]
należy do zbioru \(\mathbb{N}\)?}
{Siedem rozwiązań  w zbiorze \(\mathbb{N}\).}{Brak rozwiązań w zbiorze \(\mathbb{N}\).}{Pięć rozwiązań w zbiorze \(\mathbb{N}\).}{Więcej niż siedem rozwiązań w zbiorze  \(\mathbb{N}\).}{}{}}
\LANGcs{
\Question{
Kolik řešení má nerovnice \(\sqrt{x + 17} > x - 3,\) 
je-li \(x\in \mathbb{N}\)?}
{Právě \(7\) 
řešení.}{Nerovnice nemá v \(\mathbb{N}\) 
řešení.}{Právě \(5\) 
řešení.}{Více než \(7\) 
řešení.}{}{}}
\LANGsk{
\Question{
Koľko riešení má nerovnica 
\[ 
 \sqrt{x + 17} > x - 3
\]
 v množine \(\mathbb{N}\)?}
{Práve 7 riešení v \(\mathbb{N}\).}{Nerovnica nemá v \(\mathbb{N}\) žiadne riešenie.}{Práve päť riešení v \(\mathbb{N}\).}{Viac ako sedem riešení v \(\mathbb{N}\).}{}{}}
\LANGes{
\Question{
¿Cuántas soluciones tiene la siguiente inecuación en \(\mathbb{N}\)? 
\[ 
 \sqrt{x + 17} > x - 3
\]}
{7 soluciones}{No tiene solución en \(\mathbb{N}\).}{5 soluciones}{Más de 7 soluciones}{}{}}
\End

\Data\ID{9000024806}
\Updated{2020-12-29 09:12:03}
\Level{B}
\SubArea{Radical Equations and Inequalities}
\LANGen{
\Question{
In the following list identify the interval which is a subset of the solution set of
the following inequality. 
\[ 
 \sqrt{x^{2 } + 2x - 3} > x + 2
\]}
{\((-\infty ;-3] \)}{\(\left (-\frac{7}
{2};+\infty \right )\)}{\((1;+\infty )\)}{\((-\infty ;-2)\)}{}{}}
\LANGpl{
\Question{
Z podanych odpowiedzi wybierz przedział stanowiący podzbiór zbioru rozwiązań poniższej nierówności. 
\[ 
 \sqrt{x^{2 } + 2x - 3} > x + 2
\]}
{\((-\infty ;-3] \)}{\(\left (-\frac{7}
{2};+\infty \right )\)}{\((1;+\infty )\)}{\((-\infty ;-2)\)}{}{}}
\LANGcs{
\Question{
Je dána nerovnice \(\sqrt{x^{2 } + 2x - 3} > x + 2\).
Z následujících intervalů vyberte ty, které jsou částí množiny řešení
dané nerovnice.}
{\((-\infty ;-3\rangle \)}{\(\left (-\frac{7}
{2};+\infty \right )\)}{\((1;+\infty )\)}{\((-\infty ;-2)\)}{}{}}
\LANGsk{
\Question{
Z nasledujúcich intervalov vyberte ten, ktorý je časťou množiny riešenia danej nerovnice.
\[ 
 \sqrt{x^{2 } + 2x - 3} > x + 2
\]}
{\((-\infty ;-3] \)}{\(\left (-\frac{7}
{2};+\infty \right )\)}{\((1;+\infty )\)}{\((-\infty ;-2)\)}{}{}}
\LANGes{
\Question{
Dada la ecuación:
\[ 
 \sqrt{x^{2 } + 2x - 3} > x + 2
\]
Identifica el intervalo que es subconjunto del conjunto solución.}
{\((-\infty ;-3] \)}{\(\left (-\frac{7}
{2};+\infty \right )\)}{\((1;+\infty )\)}{\((-\infty ;-2)\)}{}{}}
\End

\Data\ID{9000024809}
\Updated{2020-07-15 03:07:34}
\Level{B}
\SubArea{Radical Equations and Inequalities}
\LANGen{
\Question{
Find the solution set of the following inequality. 
\[ 
 \sqrt{x + 3} > x - 3
\]}
{\([ -3;6)\)}{\( (1;6)\)}{\([ -3;3] \)}{\( (-\infty ;1)\cup (6;+\infty )\)}{}{}}
\LANGpl{
\Question{
Znajdź zbiór rozwiązań danej nierówności.
\[ 
 \sqrt{x + 3} > x - 3
\]}
{\([ -3;6)\)}{\( (1;6)\)}{\([ -3;3] \)}{\((-\infty ;1)\cup (6;+\infty )\)}{}{}}
\LANGcs{
\Question{
Určete množinu řešení nerovnice
\(\sqrt{
x + 3} > x - 3\).}
{\(\langle -3;6)\)}{\( (1;6)\)}{\(\langle -3;3\rangle \)}{\( (-\infty ;1)\cup (6;+\infty )\)}{}{}}
\LANGsk{
\Question{
Určte množinu riešení nerovnice.
\[ 
 \sqrt{x + 3} > x - 3
\]}
{\([ -3;6)\)}{\( (1;6)\)}{\([ -3;3] \)}{\((-\infty ;1)\cup (6;+\infty )\)}{}{}}
\LANGes{
\Question{
Calcula el conjunto de soluciones de la siguiente inecuación. 
\[ 
 \sqrt{x + 3} > x - 3
\]}
{\([ -3;6)\)}{\( (1;6)\)}{\([ -3;3] \)}{\( (-\infty ;1)\cup (6;+\infty )\)}{}{}}
\End

\Data\ID{1003177801}
\Updated{2020-07-15 03:07:21}
\Level{C}
\SubArea{Radical Equations and Inequalities}
\LANGen{
\Question{
Choose the domain of the expression \( \sqrt{|x-3|-4} \).}
{\( (-\infty;-1]\cup[7;\infty) \)}{\( (-1;3]\cup[7;\infty) \)}{\( [7;\infty) \)}{\( (-\infty;1)\cup(7;\infty) \)}{}{}}
\LANGpl{
\Question{
Wybierz dziedzinę wyrażenia \( \sqrt{|x-3|-4} \).}
{\( (-\infty;-1\rangle\cup\langle7;\infty) \)}{\( (-1;3\rangle\cup\langle7;\infty) \)}{\( \langle7;\infty) \)}{\( (-\infty;1)\cup(7;\infty) \)}{}{}}
\LANGcs{
\Question{
Určete definiční obor výrazu \( \sqrt{|x-3|-4} \).}
{\( (-\infty;-1\rangle\cup\langle7;\infty) \)}{\( (-1;3\rangle\cup\langle7;\infty) \)}{\( \langle7;\infty) \)}{\( (-\infty;1)\cup(7;\infty) \)}{}{}}
\LANGsk{
\Question{
Určte definičný obor výrazu \( \sqrt{|x-3|-4} \).}
{\( (-\infty;-1\rangle\cup\langle7;\infty) \)}{\( (-1;3\rangle\cup\langle7;\infty) \)}{\( \langle7;\infty) \)}{\( (-\infty;1)\cup(7;\infty) \)}{}{}}
\LANGes{
\Question{
Halla el dominio de la siguiente expresión: \( \sqrt{|x-3|-4} \).}
{\( (-\infty;-1]\cup[7;\infty) \)}{\( (-1;3]\cup[7;\infty) \)}{\( [7;\infty) \)}{\( (-\infty;1)\cup(7;\infty) \)}{}{}}
\End

\Data\ID{1003177803}
\Updated{2020-07-15 03:07:21}
\Level{C}
\SubArea{Radical Equations and Inequalities}
\LANGen{
\Question{
Choose the domain of the expression.
\[ \frac1{\sqrt{|3x-9|-\sqrt2}}  \]}
{\( \left(-\infty;3-\frac{\sqrt2}3\right)\cup\left(3+\frac{\sqrt2}3;\infty\right) \)}{\( \left(-\infty;-3-\frac{\sqrt2}3\right)\cup\left(3+\frac{\sqrt2}3;\infty\right) \)}{\( \left(-\infty;-3+\frac{\sqrt2}3\right)\cup\left(3+\frac{\sqrt2}3;\infty\right) \)}{\( \left(-\infty;-3-\frac{\sqrt2}3\right)\cup\left(-3+\frac{\sqrt2}3;\infty\right) \)}{}{}}
\LANGpl{
\Question{
Wybierz dziedzinę wyrażenia.
\[ \frac1{\sqrt{|3x-9|-\sqrt2}} \]}
{\( \left(-\infty;3-\frac{\sqrt2}3\right)\cup\left(3+\frac{\sqrt2}3;\infty\right) \)}{\( \left(-\infty;-3-\frac{\sqrt2}3\right)\cup\left(3+\frac{\sqrt2}3;\infty\right) \)}{\( \left(-\infty;-3+\frac{\sqrt2}3\right)\cup\left(3+\frac{\sqrt2}3;\infty\right) \)}{\( \left(-\infty;-3-\frac{\sqrt2}3\right)\cup\left(-3+\frac{\sqrt2}3;\infty\right) \)}{}{}}
\LANGcs{
\Question{
Určete definiční obor výrazu.
\[ \frac1{\sqrt{|3x-9|-\sqrt2}} \]}
{\( \left(-\infty;3-\frac{\sqrt2}3\right)\cup\left(3+\frac{\sqrt2}3;\infty\right) \)}{\( \left(-\infty;-3-\frac{\sqrt2}3\right)\cup\left(3+\frac{\sqrt2}3;\infty\right) \)}{\( \left(-\infty;-3+\frac{\sqrt2}3\right)\cup\left(3+\frac{\sqrt2}3;\infty\right) \)}{\( \left(-\infty;-3-\frac{\sqrt2}3\right)\cup\left(-3+\frac{\sqrt2}3;\infty\right) \)}{}{}}
\LANGsk{
\Question{
Určte definičný obor výrazu.
\[ \frac1{\sqrt{|3x-9|-\sqrt2}} \]}
{\( \left(-\infty;3-\frac{\sqrt2}3\right)\cup\left(3+\frac{\sqrt2}3;\infty\right) \)}{\( \left(-\infty;-3-\frac{\sqrt2}3\right)\cup\left(3+\frac{\sqrt2}3;\infty\right) \)}{\( \left(-\infty;-3+\frac{\sqrt2}3\right)\cup\left(3+\frac{\sqrt2}3;\infty\right) \)}{\( \left(-\infty;-3-\frac{\sqrt2}3\right)\cup\left(-3+\frac{\sqrt2}3;\infty\right) \)}{}{}}
\LANGes{
\Question{
Halla el dominio de la siguiente expresión:
\[ \frac1{\sqrt{|3x-9|-\sqrt2}}  \]}
{\( \left(-\infty;3-\frac{\sqrt2}3\right)\cup\left(3+\frac{\sqrt2}3;\infty\right) \)}{\( \left(-\infty;-3-\frac{\sqrt2}3\right)\cup\left(3+\frac{\sqrt2}3;\infty\right) \)}{\( \left(-\infty;-3+\frac{\sqrt2}3\right)\cup\left(3+\frac{\sqrt2}3;\infty\right) \)}{\( \left(-\infty;-3-\frac{\sqrt2}3\right)\cup\left(-3+\frac{\sqrt2}3;\infty\right) \)}{}{}}
\End

\Data\ID{1003187202}
\Updated{2020-07-15 03:07:21}
\Level{C}
\SubArea{Radical Equations and Inequalities}
\LANGen{
\Question{
Which of following equalities is true for all real numbers \( x \)?}
{\( \sqrt{(x-3)^2} = |x-3| \)}{\( |-x| = x \)}{\( |x-1|=x-1 \)}{\( \sqrt{x^2} = x \)}{}{}}
\LANGpl{
\Question{
Które z podanych równań jest prawdziwe dla wszystkich rzeczywistych liczb \( x \)?}
{\( \sqrt{(x-3)^2} = |x-3| \)}{\( |-x| = x \)}{\( |x-1|=x-1 \)}{\( \sqrt{x^2} = x \)}{}{}}
\LANGcs{
\Question{
Která z následujících rovností je pravdivá pro všechna reálná čísla  \( x \)?}
{\( \sqrt{(x-3)^2} = |x-3| \)}{\( |-x| = x \)}{\( |x-1|=x-1 \)}{\( \sqrt{x^2} = x \)}{}{}}
\LANGsk{
\Question{
Ktorá z nasledujúcich rovností je pravdivá pre všetky reálne čísla \( x \)?}
{\( \sqrt{(x-3)^2} = |x-3| \)}{\( |-x| = x \)}{\( |x-1|=x-1 \)}{\( \sqrt{x^2} = x \)}{}{}}
\LANGes{
\Question{
¿Cuál de las siguientes ecuaciones es verdadera para todos los números reales \( x \)?}
{\( \sqrt{(x-3)^2} = |x-3| \)}{\( |-x| = x \)}{\( |x-1|=x-1 \)}{\( \sqrt{x^2} = x \)}{}{}}
\End

\Data\ID{1003187203}
\Updated{2020-07-15 03:07:21}
\Level{C}
\SubArea{Radical Equations and Inequalities}
\LANGen{
\Question{
How many solutions does the equation \( \sqrt{x^2-4x+4}=|x| \) have?}
{\( 1 \)}{\( 2 \)}{\( 0 \)}{\( 3 \)}{}{}}
\LANGpl{
\Question{
Ile rozwiązań ma równanie\( \sqrt{x^2-4x+4}=|x| \)?}
{\( 1 \)}{\( 2 \)}{\( 0 \)}{\( 3 \)}{}{}}
\LANGcs{
\Question{
Kolik řešení má rovnice \( \sqrt{x^2-4x+4}=|x| \)?}
{\( 1 \)}{\( 2 \)}{\( 0 \)}{\( 3 \)}{}{}}
\LANGsk{
\Question{
Koľko riešení má rovnica \( \sqrt{x^2-4x+4}=|x| \)?}
{\( 1 \)}{\( 2 \)}{\( 0 \)}{\( 3 \)}{}{}}
\LANGes{
\Question{
¿Cuántas soluciones tiene la siguiente ecuación? \( \sqrt{x^2-4x+4}=|x| \)}
{\( 1 \)}{\( 2 \)}{\( 0 \)}{\( 3 \)}{}{}}
\End

\Data\ID{1003187204}
\Updated{2020-07-15 03:07:21}
\Level{C}
\SubArea{Radical Equations and Inequalities}
\LANGen{
\Question{
The solution set of the equation \( \sqrt{x^2-2x+1}-2|x+3|+x+7=0 \) is:}
{\( \{-7\}\cup[1;+\infty) \)}{\( \{-7\}\cup(1;+\infty) \)}{\( [1;+\infty) \)}{\( (1;+\infty) \)}{}{}}
\LANGpl{
\Question{
Zbiorem rozwiązań równania  \( \sqrt{x^2-2x+1}-2|x+3|+x+7=0 \) jest:}
{\( \{-7\}\cup\langle1;+\infty) \)}{\( \{-7\}\cup(1;+\infty) \)}{\( \langle1;+\infty) \)}{\( (1;+\infty) \)}{}{}}
\LANGcs{
\Question{
Množinou řešení rovnice \( \sqrt{x^2-2x+1}-2|x+3|+x+7=0 \) je:}
{\( \{-7\}\cup\langle1;+\infty) \)}{\( \{-7\}\cup(1;+\infty) \)}{\( \langle1;+\infty) \)}{\( (1;+\infty) \)}{}{}}
\LANGsk{
\Question{
Množinou riešení rovnice \( \sqrt{x^2-2x+1}-2|x+3|+x+7=0 \) je:}
{\( \{-7\}\cup\langle1;+\infty) \)}{\( \{-7\}\cup(1;+\infty) \)}{\( \langle1;+\infty) \)}{\( (1;+\infty) \)}{}{}}
\LANGes{
\Question{
Calcula el conjunto de soluciones de la siguiente ecuación: \( \sqrt{x^2-2x+1}-2|x+3|+x+7=0 \)}
{\( \{-7\}\cup[1;+\infty) \)}{\( \{-7\}\cup(1;+\infty) \)}{\( [1;+\infty) \)}{\( (1;+\infty) \)}{}{}}
\End

\Data\ID{9000024805}
\Updated{2020-07-15 03:07:34}
\Level{C}
\SubArea{Radical Equations and Inequalities}
\LANGen{
\Question{
A falling body dropped at a velocity
\(60\, \mathrm{m}\mathrm{s}^{-1}\). Find the
initial height \(h\),
if the relation between the velocity and the initial height
\(h\) is
\(v = \sqrt{2hg}\). Use
\(g = 10\, \mathrm{m}\mathrm{s}^{-2}\) for
acceleration of gravity.}
{The initial height is between \(150\, \mathrm{m}\) 
and \(200\, \mathrm{m}\).}{The initial height is smaller than \(100\, \mathrm{m}\).}{The initial height is between \(100\, \mathrm{m}\) 
and \(150\, \mathrm{m}\).}{The initial height is bigger than \(200\, \mathrm{m}\).}{}{}}
\LANGpl{
\Question{
Ciało spadało z prędkością 
\(60\, \mathrm{m}\mathrm{s}^{-1}\). Określ początkową wysokość \(h\),
jeśli zależność pomiędzy prędkością a początkową wysokością 
\(h\) jest równa
\(v = \sqrt{2hg}\). Użyj 
\(g = 10\, \mathrm{m}\mathrm{s}^{-2}\) jako wartości przyspieszenia ziemskiego.}
{Początkowa wysokość mieści się pomiędzy \(150\, \mathrm{m}\) 
a \(200\, \mathrm{m}\).}{Początkowa wysokość jest mniejsza niż \(100\, \mathrm{m}\).}{Początkowa wysokość mieści się pomiędzy  \(100\, \mathrm{m}\) 
i \(150\, \mathrm{m}\).}{Początkowa wysokość  jest większa niż \(200\, \mathrm{m}\).}{}{}}
\LANGcs{
\Question{
Z jaké výšky padalo těleso volným pádem, jestliže dopadlo rychlostí
\(60\, \mathrm{m}\mathrm{s}^{-1}\)?
Rychlost dopadu při volném pádu vyjadřuje vztah
\(v = \sqrt{2hg}\).
Za tíhové zrychlení dosazujte zaokrouhlenou hodnotu
\(g = 10\, \mathrm{m}\mathrm{s}^{-2}\).}
{Těleso padalo z výšky větší než
\(150\, \mathrm{m}\), ale menší
než \(200\, \mathrm{m}\).}{Těleso padalo z výšky menší než
\(100\, \mathrm{m}\).}{Těleso padalo z výšky větší než
\(100\, \mathrm{m}\), ale menší
než \(150\, \mathrm{m}\).}{Těleso padalo z výšky větší než
\(200\, \mathrm{m}\).}{}{}}
\LANGsk{
\Question{
Z akej výšky padalo teleso voľným pádom, ak dopadlo rýchlosťou
\(60\, \mathrm{m}\mathrm{s}^{-1}\)? Rýchlosť dopadu pri voľnom páde vyjadruje vzťah
\(v = \sqrt{2hg}\). Za gravitačné zrýchlenie dosadzujte zaokrúhlenú hodnotu
\(g = 10\, \mathrm{m}\mathrm{s}^{-2}\).}
{Teleso padalo z výšky väčšej ako \(150\, \mathrm{m}\),
ale menšej ako \(200\, \mathrm{m}\).}{Teleso padalo z výšky menšej ako  \(100\, \mathrm{m}\).}{Teleso padalo z výšky väčšej ako \(100\, \mathrm{m}\), 
ale menšej ako \(150\, \mathrm{m}\).}{Teleso padalo z výšky väčšej ako \(200\, \mathrm{m}\).}{}{}}
\LANGes{
\Question{
Un cuerpo cae a una velocidad de
\(60\, \mathrm{m}\mathrm{s}^{-1}\). Halla la altura inicial \(h\),
suponiendo que
\(h\) es
\(v = \sqrt{2hg}\). Utiliza
\(g = 10\, \mathrm{m}\mathrm{s}^{-2}\) para la aceleración de la gravedad.}
{La altura inicial varía desde \(150\, \mathrm{m}\) 
hasta \(200\, \mathrm{m}\).}{La altura inicial es menor de \(100\, \mathrm{m}\).}{La altura inicial varía desde \(100\, \mathrm{m}\) 
hasta \(150\, \mathrm{m}\).}{La altura inicial es mayor de \(200\, \mathrm{m}\).}{}{}}
\End

\Data\ID{9000024807}
\Updated{2020-12-29 09:12:40}
\Level{C}
\SubArea{Radical Equations and Inequalities}
\LANGen{
\Question{
A body hangs on a string of the length
\(l_{1}\). The length
\(l\) of the spring
defines the period \(T\) 
of motion by the relation 
\[ 
 T = 2\pi \sqrt{ \frac{l}
{g}},
\]
 where \(g\) 
is a standard acceleration of gravity. We have to adjust the length of the string such
that the period doubles. Find the new length of the string.}
{We elongate the string by \(3\cdot l_{1}\),
i.e. \(l_{2} = l_{1} + 3l_{1}\).}{The length doubles, i.e. \(l_{2} = 2l_{1}\).}{The new length will be half of the original length, i.e.
\(l_{2} = \frac{1} 
{2}l_1\).}{We shorten the string by \(3\cdot l_{1}\),
i.e. \(l_{2} = l_{1} - 3l_{1}\).}{}{}}
\LANGpl{
\Question{
Przedmiot wisi na sznurku o długości 
\(l_{1}\). Zachodzi zależność pomiędzy długością sznurka
\(l\) czasem ruchu  \(T\) 
opisana wzorem
\[ 
 T = 2\pi \sqrt{ \frac{l}
{g}},
\]
 gdzie \(g\) 
jest standardowym przyspieszeniem ziemskim. Jak należy dostosować długość sznurka, by czas ruchu  się podwoił?}
{Należy wydłużyć sznurek o  \(3\cdot l_{1}\),
i.e. \(l_{2} = l_{1} + 3l_{1}\).}{Należy podwoić długość tzn.  \(l_{2} = 2l_{1}\).}{Nowa długość sznurka będzie połową początkowej długości
\(l_{2} = \frac{1} 
{2}l_1\).}{Należy skrócić sznurek o  \(3\cdot l_{1}\),
i.e. \(l_{2} = l_{1} - 3l_{1}\).}{}{}}
\LANGcs{
\Question{
Těleso je zavěšeno na vlákně o délce
\(l_{1}\). Jak
musíme změnit délku vlákna, aby nově vytvořené kyvadlo kmitalo s
dvojnásobnou periodou, než kyvadlo s původní délkou? Perioda kyvadla
\(T\) závisí na
jeho délce \(l\) 
vztahem \(T = 2\pi \sqrt{ \frac{l}
{g}}\),
kde \(g\) je
tíhové zrychlení.}
{Délku zvětšíme o hodnotu \(3\cdot l_{1}\),
tj. \(l_{2} = l_{1} + 3l_{1}\).}{Délku dvakrát zvětšíme, tj.
\(l_{2} = 2l_{1}\).}{Délku dvakrát zmenšíme, tj. \(l_{2} = \frac{1} 
{2}l_1\).}{Délku zmenšíme o hodnotu \(3\cdot l_{1}\),
tj. \(l_{2} = l_{1} - 3l_{1}\).}{}{}}
\LANGsk{
\Question{
Teleso je zavesené na vlákne s dĺžkou
\(l_{1}\). Ako musíme zmeniť dĺžku vlákna, aby novovytvorené kyvadlo kmitalo s dvojnásobnou periódou než kyvadlo s pôvodnou dĺžkou?
  Perióda kyvadla \(T\) závisí na jeho dĺžke  vzťahom
 \[ 
 T = 2\pi \sqrt{ \frac{l}
{g}},
\]
 kde \(g\) 
je gravitačné zrýchlenie.}
{Dĺžku zväčšíme o hodnotu \(3\cdot l_{1}\),
t.j. \(l_{2} = l_{1} + 3l_{1}\).}{Dĺžku dvakrát zväčšíme, t.j. \(l_{2} = 2l_{1}\).}{Dĺžku dvakrát zmenšíme, t.j.
\(l_{2} = \frac{1} 
{2}l_1\).}{Dĺžku zmenšíme o hodnotu  \(3\cdot l_{1}\),
t.j. \(l_{2} = l_{1} - 3l_{1}\).}{}{}}
\LANGes{
\Question{
Un cuerpo está colgado de un hilo de longitud
\(l_{1}\). La longitud
\(l\) del hilo
define el período \(T\) 
de movimiento por la relación 
\[ 
 T = 2\pi \sqrt{ \frac{l}
{g}},
\]
 donde \(g\) 
es la aceleración estándar de la gravedad. Tenemos que ajustar la longitud del hilo para que el período se duplique. Halla la nueva longitud del hilo.}
{Alargamos el hilo por \(3\cdot l_{1}\),
e.d. \(l_{2} = l_{1} + 3l_{1}\).}{La longitud la duplicamos, e.d. \(l_{2} = 2l_{1}\).}{La nueva longitud será la mitad de la longitud original, e.d.
\(l_{2} = \frac{1} 
{2}l_1\).}{Acortamos el hilo por \(3\cdot l_{1}\),
e.d. \(l_{2} = l_{1} - 3l_{1}\).}{}{}}
\End

\Data\ID{9000024808}
\Updated{2020-12-29 09:12:01}
\Level{C}
\SubArea{Radical Equations and Inequalities}
\LANGen{
\Question{
In the following list identify a true statement referring to the following equation.
\[ 
 \sqrt{4x^{2 } - \sqrt{8x + 5}} = 2x + 1
\]}
{The equation has a unique solution, this solution is a negative number.}{The equation has two solutions, both solutions have an opposite sign.}{The equation has a unique solution, this solution is a positive number.}{The equation does not have a solution.}{}{}}
\LANGpl{
\Question{
Które z poniższych zdań dotyczących poniższego równania jest prawdziwe?
\[ 
 \sqrt{4x^{2 } - \sqrt{8x + 5}} = 2x + 1
\]}
{Równanie ma tylko jedno rozwiązanie i jest nim liczba ujemna.}{Równanie ma dwa rozwiązania przeciwnych znaków.}{Równanie ma tylko jedno rozwiązanie i jest nim liczba dodatnia.}{Równanie nie ma rozwiązania.}{}{}}
\LANGcs{
\Question{
Je dána rovnice \(\sqrt{4x^{2 } - \sqrt{8x + 5}} = 2x + 1\).
Vyberte pravdivé tvrzení o kořenech této rovnice.}
{Rovnice má právě jeden záporný kořen.}{Rovnice má právě \(2\) 
kořeny, které se liší znaménkem.}{Rovnice má právě jeden kladný kořen.}{Rovnice nemá žádný kořen.}{}{}}
\LANGsk{
\Question{
Je daná rovnica.
\[ 
 \sqrt{4x^{2 } - \sqrt{8x + 5}} = 2x + 1
\] Vyberte pravdivé tvrdenie o koreňoch tejto rovnice.}
{Rovnica má práve jeden záporný koreň.}{Rovnica má práve dva korene, ktoré sa líšia znamienkom.}{Rovnica má práve jeden kladný koreň.}{Rovnica nemá žiadny koreň.}{}{}}
\LANGes{
\Question{
Dada la ecuación:
\[ 
 \sqrt{4x^{2 } - \sqrt{8x + 5}} = 2x + 1
\]
Identifica la proposición lógica que hace referencia a la ecuación.}
{La ecuación tiene solo una solución, es un número negativo.}{La ecuación tiene dos soluciones, ambas soluciones tienen signo contrario.}{La ecuación tiene solo una solución, es un número positivo.}{La ecuación no tiene soluciones.}{}{}}
\End

\Data\ID{9000024810}
\Updated{2020-07-15 03:07:34}
\Level{C}
\SubArea{Radical Equations and Inequalities}
\LANGen{
\Question{
Find the solution set of the following inequality. 
\[ 
 \sqrt{|x|} > x
\]}
{\(K = (-\infty ;0)\cup (0;1)\)}{\(K = (-\infty ;1)\)}{\(K =\mathbb{R} ^{+}\)}{\(K = (0;1)\)}{}{}}
\LANGpl{
\Question{
Znajdź zbiór rozwiązań danej nierówności:
\[ 
 \sqrt{|x|} > x
\]}
{\(K = (-\infty ;0)\cup (0;1)\)}{\(K = (-\infty ;1)\)}{\(K =\mathbb{R} ^{+}\)}{\(K = (0;1)\)}{}{}}
\LANGcs{
\Question{
Určete množinu řešení nerovnice
\(\sqrt{
|x|} > x\).}
{\(K = (-\infty ;0)\cup (0;1)\)}{\(K = (-\infty ;1)\)}{\(K =\mathbb{R} ^{+}\)}{\(K = (0;1)\)}{}{}}
\LANGsk{
\Question{
Určte množinu riešení danej nerovnice. 
\[ 
 \sqrt{|x|} > x
\]}
{\(K = (-\infty ;0)\cup (0;1)\)}{\(K = (-\infty ;1)\)}{\(K =\mathbb{R} ^{+}\)}{\(K = (0;1)\)}{}{}}
\LANGes{
\Question{
Calcula el conjunto de soluciones de la siguiente inecuación.
\[ 
 \sqrt{|x|} > x
\]}
{\(K = (-\infty ;0)\cup (0;1)\)}{\(K = (-\infty ;1)\)}{\(K =\mathbb{R} ^{+}\)}{\(K = (0;1)\)}{}{}}
\End
